\chapter{Fuzzy Frames, Fuzzy Locales, and the Fuzzy Sierpiński Classifier}
\label{ch:fuzzy-frames-locales}

In the previous chapter we constructed fuzzy predicates as enriched presheaves. For a
small \(\mathbb V\)-category \(X\), the fuzzy presheaf category was
\[
    P_{\mathbb V}(X)
    =
    [X^{\op},\underline{\mathbb V}],
\]
and, after passing to the posetal specialization determined by the fuzzy truth-value
object
\[
    \Lambda=(L,\leq,\meet,\join,\tensor,\multimap,\bot_L,\top_L),
\]
the corresponding fuzzy predicate object was
\[
    P_\Lambda(X)
    =
    [X^{\op},\Lambda].
\]
This chapter gives these presheaf objects their localic interpretation.

There are two structures on \(P_\Lambda(X)\) which should be kept distinct. First,
\(P_\Lambda(X)\) is a separated \(\Lambda\)-preorder, with hom-values
\[
    P_\Lambda(X)(\phi,\psi)
    =
    \bigmeet_{x\in X}(\phi(x)\multimap\psi(x)).
\]
Second, its induced crisp order is a frame, with joins and meets computed pointwise.
In this chapter, when we call \(P_\Lambda(X)\) a frame, we mean this induced crisp
frame of predicates, equipped with the pointwise \(\Lambda\)-scalar action
\[
    (r\tensor\phi)(x)=r\tensor\phi(x).
\]
The enriched hom-values remain present, but the frame order is the induced crisp
order.

The guiding idea is that fuzzy frames are obtained by localizing fuzzy presheaves. At
the categorical level this is expressed by enriched left-exact reflective
localizations of
\[
    P_{\mathbb V}(X).
\]
At the posetal level this becomes a scalar-compatible nucleus on
\[
    P_\Lambda(X).
\]
We prove below that these two descriptions agree in the posetal specialization:
left-exact \(\Lambda\)-enriched reflective localizations of \(P_\Lambda(X)\)
correspond, up to equivalence over \(P_\Lambda(X)\), to fuzzy nuclei on
\(P_\Lambda(X)\). Fuzzy locales are then defined by formal duality.

The fuzzy Sierpiński classifier is treated by its classifying universal property. This
is an important point. In the crisp case the Sierpiński frame can be realized as the
presheaf frame on the two-element preorder
\[
    U\leq 1.
\]
For general fuzzy truth-value bases this raw presheaf construction does not, in
general, have the required free-open property. In particular, for non-idempotent
tensors such as the {\L}ukasiewicz and product tensors, the presheaf frame on the
two-point subterminal shape need not classify arbitrary fuzzy opens by
finite-meet-preserving scalar frame homomorphisms.

Accordingly, the fuzzy Sierpiński frame is defined in this chapter as the free fuzzy
frame on one fuzzy open:
\[
    \Omega_\Lambda=F_{\op}^{\Lambda}(1),
\]
with generic fuzzy open
\[
    \sigma_\Lambda\in\Omega_\Lambda.
\]
Its dual fuzzy locale
\[
    \mathbb S_\Lambda
\]
classifies fuzzy opens by the universal property
\[
    \FLoc_\Lambda(\mathcal X,\mathbb S_\Lambda)
    \cong
    |\Opens_\Lambda(\mathcal X)|.
\]
Thus the Sierpiński object is not identified with a raw presheaf object unless an
additional presentation theorem has been proved for the chosen base.

At the enriched-cosmos level, the analogous construction is bicategorical. The free
fuzzy frame on an \(S\)-indexed family of cosmos-level fuzzy opens represents a
hom-category of models:
\[
    \mathbf{FFrm}_{\mathbb V}
    \bigl(F_{\op}(S),\mathcal A\bigr)
    \simeq
    \operatorname{Sub}_{\mathcal A}(1_{\mathcal A})^S.
\]
Equivalently, if \(S_{\mathrm{disc}}\) is the discrete category on \(S\), the
right-hand side is
\[
    \mathbf{Cat}
    \bigl(
        S_{\mathrm{disc}},
        \operatorname{Sub}_{\mathcal A}(1_{\mathcal A})
    \bigr).
\]
After passing to the posetal reflection, these hom-categories become ordinary
hom-sets and the corresponding classification becomes a strict set-level universal
property.

All large categories and collections of subobjects are taken relative to the fixed
universe convention of the previous chapter. If \(\mathcal A\) is an ordinary or
enriched category with terminal object, we write
\[
    \operatorname{Sub}_{\mathcal A}(1_{\mathcal A})
\]
for the ordinary category of subobjects of \(1_{\mathcal A}\). Its objects are
monomorphisms
\[
    U\hookrightarrow 1_{\mathcal A},
\]
and its morphisms are commuting triangles over \(1_{\mathcal A}\). When only
isomorphism classes are needed, we write
\[
    \operatorname{Iso}\operatorname{Sub}_{\mathcal A}(1_{\mathcal A})
    :=
    \operatorname{Ob}\operatorname{Sub}_{\mathcal A}(1_{\mathcal A})/{\cong}.
\]
In the posetal reflection, \(\operatorname{Sub}_{A}(\top_A)\) is equivalent to the
thin category underlying the frame \(A\), and its set of objects is simply \(|A|\).

\paragraph{Finite-limit convention for this chapter.}
We fix once and for all a class \(\mathcal F\) of finite weighted limits in
\(\mathbb V\)-categories. The class \(\mathcal F\) contains the terminal-object and
binary-product weights and is the finite-limit class used in all terms such as
``left exact'', ``finite-limit preserving'', and
\(\operatorname{Lex}_{\mathbb V}\). Thus
\[
    \operatorname{Lex}_{\mathbb V}(\mathcal C,\mathcal A)
\]
denotes the ordinary category of \(\mathbb V\)-functors
\[
    \mathcal C\to\mathcal A
\]
preserving the chosen \(\mathcal F\)-limits, with enriched natural transformations.

In the posetal results below, the only finite limits used explicitly are the
terminal element and binary meets of the underlying crisp frame. Thus left exactness
in the posetal comparison theorem means preservation of \(\top\) and binary meets.

The overall progression is:
\[
\begin{tikzcd}[column sep=large]
\text{fuzzy preorders}
    \arrow[r, "{P_\Lambda}"']
&
\text{fuzzy presheaves}
    \arrow[r, "\text{localization}"']
&
\text{fuzzy frames}
    \arrow[r, "\op"']
&
\text{fuzzy locales}
    \arrow[r, "{\FLoc_\Lambda(-,\mathbb S_\Lambda)}"']
&
\text{fuzzy opens.}
\end{tikzcd}
\]

The notation \(P_{\mathbb V}(X)\) and \(P_\Lambda(X)\) is retained from the previous
chapter. The new notation introduced in this chapter is:
\[
    \FFrm_\Lambda
\]
for the category of fuzzy frames over \(\Lambda\),
\[
    \FLoc_\Lambda
\]
for the category of fuzzy locales over \(\Lambda\),
\[
    \Omega_\Lambda
\]
for the fuzzy Sierpiński frame,
\[
    \sigma_\Lambda\in\Omega_\Lambda
\]
for its generic fuzzy open, and
\[
    \mathbb S_\Lambda
\]
for the corresponding fuzzy Sierpiński locale.

\section{Fuzzy presheaves recalled}
\label{sec:fuzzy-presheaves-recalled}

We begin by recalling the presheaf constructions from the previous chapter. The
definition--theorem--proof development of these objects has already been given there.

For a small \(\mathbb V\)-category \(X\), the fuzzy presheaf category is
\[
    P_{\mathbb V}(X)
    =
    [X^{\op},\underline{\mathbb V}].
\]
This is generally a locally small, possibly large \(\mathbb V\)-category. Its objects
are fuzzy presheaves on \(X\). The Yoneda embedding is the \(\mathbb V\)-functor
\[
    \yon_X:X\to P_{\mathbb V}(X),
    \qquad
    x\longmapsto X(-,x).
\]
By the enriched Yoneda lemma,
\[
    P_{\mathbb V}(X)(\yon_Xx,\Phi)
    \cong
    \Phi x
\]
for every fuzzy presheaf \(\Phi\in P_{\mathbb V}(X)\). Equivalently, the
variance-correct Yoneda comparison says that the composite
\[
\begin{tikzcd}[column sep=large, row sep=large]
X^{\op}
    \arrow[r, "\yon_X^{\op}"]
    \arrow[dr, "\Phi"']
&
P_{\mathbb V}(X)^{\op}
    \arrow[d, "{P_{\mathbb V}(X)(-,\Phi)}"]
\\
&
\underline{\mathbb V}
\end{tikzcd}
\]
is isomorphic to \(\Phi\). These are the standard enriched presheaf and Yoneda
constructions \cite{Kelly1982BasicConcepts}.

The density formula says that every fuzzy presheaf is a weighted colimit of
representables:
\[
    \Phi
    \cong
    \Phi\star\yon_X
    \cong
    \int^{x\in X}\Phi x\tensor\yon_Xx.
\]
Thus the values of \(\Phi\) give the coefficients with which the principal presheaves
\(\yon_Xx\) are joined to reconstruct \(\Phi\). This is the usual enriched density
formula for the Yoneda embedding \cite{Kelly1982BasicConcepts}.

After passing to the posetal specialization determined by
\[
    \Lambda=(L,\leq,\meet,\join,\tensor,\multimap,\bot_L,\top_L),
\]
a fuzzy preorder \(X\) has fuzzy predicate object
\[
    P_\Lambda(X)
    =
    [X^{\op},\Lambda].
\]
An element of \(P_\Lambda(X)\) is a map
\[
    \phi:X\to L
\]
satisfying
\[
    X(x',x)\tensor\phi(x)\leq \phi(x')
\]
for all \(x,x'\in X\). Joins, meets, and scalar multiplication in the induced crisp
frame are computed pointwise:
\[
    \left(\bigjoin_i\phi_i\right)(x)
    =
    \bigjoin_i\phi_i(x),
\]
\[
    \left(\bigmeet_i\phi_i\right)(x)
    =
    \bigmeet_i\phi_i(x),
\]
and
\[
    (r\tensor\phi)(x)
    =
    r\tensor\phi(x).
\]
This is the form in which the presheaf construction will be localized below.

\section{Localizing fuzzy presheaves}
\label{sec:localizing-fuzzy-presheaves}

We now perform the localic operation. Instead of postulating an arbitrary algebraic
notion of fuzzy frame, we obtain fuzzy frames by localizing fuzzy presheaves.

At the categorical level, this means a left-exact enriched localization of
\[
    P_{\mathbb V}(X).
\]
At the posetal level, this becomes a scalar-compatible nucleus on
\[
    P_\Lambda(X).
\]
The fixed points of such a nucleus form the corresponding fuzzy frame.

\subsection{Categorical form}
\label{subsec:localizing-presheaves-categorical-form}

Let \(X\) be a small \(\mathbb V\)-category. A localization of the fuzzy presheaf
category \(P_{\mathbb V}(X)\) is given by a \(\mathbb V\)-adjunction
\[
    a\dashv i:
    P_{\mathbb V}(X)\rightleftarrows \mathcal A
\]
with \(i\) fully faithful. The associated idempotent monad is
\[
    J\eqdef ia.
\]
We require this localization to be left exact with respect to the chosen
\(\mathcal F\)-limits.

\begin{definition}[Categorical fuzzy-frame presentation]
\label{def:fuzzy-frame-presentation-categorical}
A \textbf{categorical fuzzy-frame presentation} is a \(\mathbb V\)-adjunction
\[
    a\dashv i:
    P_{\mathbb V}(X)\rightleftarrows \mathcal A
\]
between locally small, possibly large \(\mathbb V\)-categories such that:
\begin{enumerate}
    \item \(i\) is fully faithful as a \(\mathbb V\)-functor;
    \item \(\mathcal A\) has the chosen finite weighted limits;
    \item \(a\) preserves the chosen finite weighted limits.
\end{enumerate}
The fuzzy frame presented by this localization has category of opens \(\mathcal A\).
\end{definition}

The localization is displayed by
\[
\begin{tikzcd}[column sep=large, row sep=large]
P_{\mathbb V}(X)
    \arrow[r, shift left=1.1ex, "a"]
&
\mathcal A
    \arrow[l, shift left=1.1ex, "i"]
\arrow[phantom, from=1-1, to=1-2, "\dashv" description]
\end{tikzcd}
\]
and the associated idempotent monad is
\[
\begin{tikzcd}[column sep=large]
P_{\mathbb V}(X)
    \arrow[r, "a"]
&
\mathcal A
    \arrow[r, "i"]
&
P_{\mathbb V}(X).
\end{tikzcd}
\]
Weighted limits and colimits are understood in the standard enriched sense
\cite{Kelly1982BasicConcepts}.

\begin{definition}[Local object]
\label{def:local-object-categorical}
Let
\[
    J=ia:P_{\mathbb V}(X)\to P_{\mathbb V}(X)
\]
be the idempotent monad associated to a categorical fuzzy-frame presentation. A fuzzy
presheaf \(\Phi\in P_{\mathbb V}(X)\) is called \textbf{\(J\)-local} if the unit
\[
    \Phi\to J\Phi
\]
is an isomorphism.
\end{definition}

\begin{proposition}[Local objects]
\label{prop:local-objects-categorical}
Let
\[
    a\dashv i:
    P_{\mathbb V}(X)\rightleftarrows \mathcal A
\]
be a categorical fuzzy-frame presentation, and let
\[
    J=ia.
\]
Then the full subcategory of \(J\)-local objects in \(P_{\mathbb V}(X)\) is equivalent
to \(\mathcal A\).
\end{proposition}

\begin{proof}
Since \(i\) is fully faithful, the counit
\[
    aiA\to A
\]
is an isomorphism for every \(A\in\mathcal A\). Hence
\[
    J(iA)=iaiA\cong iA,
\]
so every object in the image of \(i\) is \(J\)-local.

Conversely, if \(\Phi\) is \(J\)-local, then
\[
    \Phi\cong J\Phi=ia\Phi,
\]
so \(\Phi\) is isomorphic to an object in the image of \(i\). Therefore the local
objects are precisely the objects of \(\mathcal A\), up to equivalence.
\end{proof}

The local-object subcategory sits inside the presheaf category as
\[
\begin{tikzcd}[column sep=large, row sep=large]
\mathcal A
    \arrow[r, hook, "i"]
    \arrow[dr, "\simeq"']
&
P_{\mathbb V}(X)
\\
&
\Fix(J)
    \arrow[u, hook']
\end{tikzcd}
\]
where \(\Fix(J)\) denotes the full subcategory of \(J\)-local objects.

\begin{proposition}[Closed colimits]
\label{prop:closed-colimits-categorical}
Let
\[
    a\dashv i:
    P_{\mathbb V}(X)\rightleftarrows \mathcal A
\]
be a categorical fuzzy-frame presentation. Let \(K\) be small, let
\[
    W:K^{\op}\to\underline{\mathbb V}
\]
be a weight, and let
\[
    D:K\to\mathcal A
\]
be a diagram. Whenever the weighted colimit \(W\star iD\) exists in
\(P_{\mathbb V}(X)\), the weighted colimit of \(D\) exists in \(\mathcal A\) and is
computed by
\[
    W\star_{\mathcal A}D
    \cong
    a(W\star iD).
\]
\end{proposition}

\begin{proof}
For \(A\in\mathcal A\), the enriched adjunction and the universal property of the
weighted colimit give
\[
\begin{aligned}
\mathcal A(a(W\star iD),A)
&\cong
P_{\mathbb V}(X)(W\star iD,iA)
\\
&\cong
[K^{\op},\underline{\mathbb V}]
\bigl(
W,
P_{\mathbb V}(X)(iD-,iA)
\bigr)
\\
&\cong
[K^{\op},\underline{\mathbb V}]
\bigl(
W,
\mathcal A(D-,A)
\bigr).
\end{aligned}
\]
This is the enriched universal property of the weighted colimit of \(D\) in
\(\mathcal A\).
\end{proof}

\begin{proposition}[Finite limits created by the inclusion]
\label{prop:finite-limits-created-categorical}
Let
\[
    a\dashv i:
    P_{\mathbb V}(X)\rightleftarrows \mathcal A
\]
be a categorical fuzzy-frame presentation. Let \(K\) and \(W\) belong to the chosen
finite-limit class \(\mathcal F\), and let
\[
    D:K\to\mathcal A
\]
be a diagram. Then the corresponding finite weighted limits in \(\mathcal A\) are
created by the inclusion
\[
    i:\mathcal A\to P_{\mathbb V}(X).
\]
\end{proposition}

\begin{proof}
Let
\[
    M=\{W,iD\}
\]
be the corresponding finite weighted limit in \(P_{\mathbb V}(X)\). Since \(a\)
preserves the chosen finite weighted limits,
\[
    aM
    \cong
    a\{W,iD\}
    \cong
    \{W,aiD\}
    \cong
    \{W,D\}
\]
in \(\mathcal A\).

Since \(i\) is a right adjoint, it preserves limits. Hence \(iaM\) is also a limit of
\(iD\) in \(P_{\mathbb V}(X)\). But \(M\) is already the limit of \(iD\). The unit
\[
    M\to iaM
\]
is the unique comparison map between two limiting cones over \(iD\). Since both
objects satisfy the same finite-limit universal property in \(P_{\mathbb V}(X)\),
this comparison is an isomorphism. Thus \(M\) is local, and the chosen finite limits
in \(\mathcal A\) are created by \(i\).
\end{proof}

Thus, categorically, localization has the expected frame-theoretic behavior: colimits
are formed in the ambient presheaf category and then closed, while finite limits are
inherited by the local objects.

\subsection{Posetal reflection and fuzzy nuclei}
\label{subsec:localizing-presheaves-posetal-reflection}

We now pass to the posetal specialization. The categorical localization above becomes
a closure operation on
\[
    P_\Lambda(X),
\]
compatible with finite meets and with the scalar action inherited from \(\Lambda\).
The first three axioms below are the usual nucleus axioms from frame theory
\cite{Johnstone1982}. The final axiom is the additional compatibility with fuzzy
truth-value rescaling.

\begin{definition}[Fuzzy nucleus]
\label{def:fuzzy-nucleus}
Let \(X\) be a fuzzy preorder. A \textbf{fuzzy nucleus} on \(P_\Lambda(X)\) is a map
\[
    j:P_\Lambda(X)\to P_\Lambda(X)
\]
such that, for all \(\phi,\psi\in P_\Lambda(X)\) and all \(r\in L\),
\[
    \phi\leq j(\phi),
\]
\[
    j(j(\phi))=j(\phi),
\]
\[
    j(\phi\meet\psi)=j(\phi)\meet j(\psi),
\]
and
\[
    j(r\tensor j(\phi))=j(r\tensor\phi).
\]
\end{definition}

The first three equations are equivalent to the usual nucleus axioms on the frame
\(P_\Lambda(X)\); monotonicity follows from finite-meet preservation and is recorded
below. The final equation says that localization commutes with fuzzy truth-value
rescaling after closure. This scalar-compatible form is close in spirit to nucleus
constructions in quantale theory \cite{Rosenthal1990}.

\begin{lemma}[Fuzzy nuclei are monotone]
\label{lem:fuzzy-nuclei-monotone}
Every fuzzy nucleus is monotone.
\end{lemma}

\begin{proof}
If \(\phi\leq\psi\), then
\[
    \phi=\phi\meet\psi.
\]
Therefore
\[
    j(\phi)
    =
    j(\phi\meet\psi)
    =
    j(\phi)\meet j(\psi)
    \leq
    j(\psi).
\]
\end{proof}

\begin{lemma}[Fuzzy nuclei preserve top]
\label{lem:fuzzy-nuclei-preserve-top}
Every fuzzy nucleus satisfies
\[
    j(\top_{P_\Lambda(X)})=\top_{P_\Lambda(X)}.
\]
\end{lemma}

\begin{proof}
Inflationarity gives
\[
    \top_{P_\Lambda(X)}\leq j(\top_{P_\Lambda(X)}).
\]
Since \(\top_{P_\Lambda(X)}\) is the greatest element of \(P_\Lambda(X)\), this
forces
\[
    j(\top_{P_\Lambda(X)})=\top_{P_\Lambda(X)}.
\]
\end{proof}

\begin{lemma}[Scalars as enriched tensors]
\label{lem:scalars-as-enriched-tensors}
Let \(P=P_\Lambda(X)\). For every \(r\in L\), the pointwise scalar operation
\[
    r\tensor(-):P\to P
\]
is the tensor of the \(\Lambda\)-category \(P\) by the scalar \(r\). Explicitly, for
all \(\phi,\psi\in P\),
\[
    P(r\tensor\phi,\psi)
    =
    r\multimap P(\phi,\psi).
\]
\end{lemma}

\begin{proof}
By the scalar-closure result for fuzzy predicates from the previous chapter,
\(r\tensor\phi\) is again an element of \(P\). By definition,
\[
    P(r\tensor\phi,\psi)
    =
    \bigmeet_{x\in X}
    \bigl((r\tensor\phi(x))\multimap\psi(x)\bigr).
\]
For \(q\in L\), residuation gives
\[
\begin{aligned}
q\leq P(r\tensor\phi,\psi)
&\Longleftrightarrow
q\tensor r\tensor\phi(x)\leq\psi(x)
\quad\text{for all }x
\\
&\Longleftrightarrow
q\tensor r\leq \phi(x)\multimap\psi(x)
\quad\text{for all }x
\\
&\Longleftrightarrow
q\tensor r\leq
\bigmeet_{x\in X}(\phi(x)\multimap\psi(x))
\\
&\Longleftrightarrow
q\tensor r\leq P(\phi,\psi)
\\
&\Longleftrightarrow
q\leq r\multimap P(\phi,\psi).
\end{aligned}
\]
Hence
\[
    P(r\tensor\phi,\psi)
    =
    r\multimap P(\phi,\psi).
\]
\end{proof}

\begin{lemma}[Tensors in a reflective localization]
\label{lem:tensors-in-reflective-localization}
Let
\[
    a\dashv i:
    P_\Lambda(X)\rightleftarrows\mathcal A
\]
be a \(\Lambda\)-enriched reflective localization with \(i\) fully faithful. Then
\(\mathcal A\) is tensored over \(\Lambda\). For \(r\in L\) and \(A\in\mathcal A\),
the tensor is computed by
\[
    r\tensor_{\mathcal A}A
    \cong
    a(r\tensor iA).
\]
Moreover,
\[
    a(r\tensor\phi)
    \cong
    r\tensor_{\mathcal A}a\phi
\]
for all \(\phi\in P_\Lambda(X)\).
\end{lemma}

\begin{proof}
For \(B\in\mathcal A\), using the enriched adjunction and Lemma
\ref{lem:scalars-as-enriched-tensors}, we compute
\[
\begin{aligned}
\mathcal A(a(r\tensor iA),B)
&\cong
P_\Lambda(X)(r\tensor iA,iB)
\\
&\cong
r\multimap P_\Lambda(X)(iA,iB)
\\
&\cong
r\multimap \mathcal A(A,B).
\end{aligned}
\]
Thus \(a(r\tensor iA)\) satisfies the universal property of
\(r\tensor_{\mathcal A}A\).

Taking \(A=a\phi\) and using the counit \(aia\phi\cong a\phi\), we have
\[
    r\tensor_{\mathcal A}a\phi
    \cong
    a(r\tensor ia\phi).
\]
Also, since \(a\) is an enriched left adjoint, it preserves tensors. Hence
\[
    a(r\tensor\phi)
    \cong
    r\tensor_{\mathcal A}a\phi.
\]
\end{proof}

\begin{theorem}[Posetal comparison for localizations and fuzzy nuclei]
\label{thm:posetal-comparison-localizations-nuclei}
Let \(X\) be a fuzzy preorder, and write
\[
    P=P_\Lambda(X).
\]
There is a correspondence, up to equivalence of reflective localizations over \(P\),
between:
\begin{enumerate}[label=\textup{(\arabic*)}]
    \item left-exact \(\Lambda\)-enriched reflective localizations
    \[
        a\dashv i:
        P\rightleftarrows\mathcal A
    \]
    with \(i\) fully faithful, where left exactness in the posetal reflection means
    preservation of the terminal element and binary meets in the underlying frame;

    \item fuzzy nuclei
    \[
        j:P\to P.
    \]
\end{enumerate}
Under this correspondence, a localization determines the fuzzy nucleus
\[
    j=ia,
\]
and a fuzzy nucleus \(j\) determines the reflective localization
\[
    q_j\dashv i_j:
    P\rightleftarrows P_j,
\]
where
\[
    P_j=\{\phi\in P\mid j(\phi)=\phi\},
    \qquad
    q_j(\phi)=j(\phi),
\]
and \(i_j:P_j\hookrightarrow P\) is the inclusion. Moreover,
\[
    \mathcal A\simeq P_j
\]
over \(P\).

In the posetal specialization, the underlying crisp order of \(P\) is a genuine
poset. Since \(P\) is separated as a \(\Lambda\)-preorder, any ordinary isomorphism
between objects of \(P\) is an equality. Thus, after replacing a localization by its
equivalent full subpreorder of local objects, the idempotent monad \(ia\) may be
regarded strictly as a closure operator on the underlying predicate frame.
\end{theorem}

\begin{proof}
First let
\[
    a\dashv i:
    P\rightleftarrows\mathcal A
\]
be a left-exact \(\Lambda\)-enriched reflective localization with \(i\) fully
faithful, and define
\[
    j=ia.
\]
The unit of the adjunction gives
\[
    \phi\leq ia\phi=j(\phi),
\]
so \(j\) is inflationary. Since \(i\) is fully faithful, the counit
\[
    aiA\to A
\]
is an isomorphism for every \(A\in\mathcal A\). Hence
\[
    j(j(\phi))
    =
    iaia\phi
    \cong
    ia\phi
    =
    j(\phi).
\]
Since \(P\) is separated, this isomorphism may be read as the strict equality
\[
    j(j(\phi))=j(\phi).
\]
Thus \(j\) is idempotent.

Because \(a\) preserves the finite meets of the underlying frame and \(i\), as a
right adjoint, preserves them, we have
\[
\begin{aligned}
j(\phi\meet\psi)
&=
ia(\phi\meet\psi)
\\
&\cong
i(a\phi\meet a\psi)
\\
&\cong
ia\phi\meet ia\psi
\\
&=
j(\phi)\meet j(\psi).
\end{aligned}
\]
Again using separatedness of \(P\), this comparison is a strict equality:
\[
    j(\phi\meet\psi)=j(\phi)\meet j(\psi).
\]
Similarly,
\[
    j(\top_P)=\top_P.
\]

It remains to prove scalar compatibility. By Lemma
\ref{lem:tensors-in-reflective-localization},
\[
    a(r\tensor\phi)
    \cong
    r\tensor_{\mathcal A}a\phi
    \cong
    a(r\tensor ia\phi).
\]
Therefore
\[
\begin{aligned}
j(r\tensor j(\phi))
&=
ia(r\tensor ia\phi)
\\
&\cong
ia(r\tensor\phi)
\\
&=
j(r\tensor\phi).
\end{aligned}
\]
Using separatedness once more, this gives the strict equality
\[
    j(r\tensor j(\phi))=j(r\tensor\phi).
\]
Thus \(j\) is a fuzzy nucleus.

Conversely, let \(j:P\to P\) be a fuzzy nucleus, and let
\[
    P_j=\{\phi\in P\mid j(\phi)=\phi\}
\]
be the full \(\Lambda\)-subpreorder of fixed points. Define
\[
    q_j:P\to P_j,
    \qquad
    q_j(\phi)=j(\phi),
\]
and let
\[
    i_j:P_j\hookrightarrow P
\]
be the inclusion.

First \(q_j\) is a \(\Lambda\)-functor. Let \(r\leq P(\phi,\psi)\). Equivalently,
\[
    r\tensor\phi\leq\psi.
\]
By monotonicity of \(j\),
\[
    j(r\tensor\phi)\leq j(\psi).
\]
By inflationarity applied to \(r\tensor j(\phi)\), followed by scalar compatibility,
\[
    r\tensor j(\phi)
    \leq
    j(r\tensor j(\phi))
    =
    j(r\tensor\phi).
\]
Hence
\[
    r\tensor j(\phi)\leq j(\psi),
\]
and so
\[
    r\leq P(j\phi,j\psi).
\]
Therefore
\[
    P(\phi,\psi)\leq P_j(q_j\phi,q_j\psi),
\]
so \(q_j\) is a \(\Lambda\)-functor.

We now prove that
\[
    q_j\dashv i_j.
\]
Let \(\psi\in P_j\). Since \(P_j\) is full,
\[
    P_j(q_j\phi,\psi)=P(j\phi,\psi).
\]
We prove
\[
    P(j\phi,\psi)=P(\phi,\psi).
\]
For \(r\in L\), if
\[
    r\leq P(j\phi,\psi),
\]
then
\[
    r\tensor j\phi\leq\psi.
\]
Since \(\phi\leq j\phi\), this implies
\[
    r\tensor\phi\leq\psi,
\]
so
\[
    r\leq P(\phi,\psi).
\]
Conversely, if
\[
    r\tensor\phi\leq\psi,
\]
then
\[
    j(r\tensor\phi)\leq j(\psi)=\psi.
\]
By inflationarity and scalar compatibility,
\[
    r\tensor j\phi
    \leq
    j(r\tensor j\phi)
    =
    j(r\tensor\phi)
    \leq
    \psi.
\]
Thus
\[
    r\leq P(j\phi,\psi).
\]
Therefore
\[
    P(j\phi,\psi)=P(\phi,\psi),
\]
and hence \(q_j\dashv i_j\) as a \(\Lambda\)-adjunction.

The inclusion \(i_j\) is fully faithful by definition. The reflector \(q_j\) preserves
terminal element and binary meets because
\[
    j(\top_P)=\top_P
\]
and
\[
    j(\phi\meet\psi)=j(\phi)\meet j(\psi).
\]
Thus
\[
    q_j\dashv i_j:
    P\rightleftarrows P_j
\]
is a left-exact \(\Lambda\)-enriched reflective localization.

Finally, starting from a localization and applying the construction above gives the
fixed-point subcategory of the idempotent monad \(j=ia\), which is equivalent to
\(\mathcal A\) by Proposition~\ref{prop:local-objects-categorical}. Starting from a
fuzzy nucleus \(j\), the induced monad is exactly \(i_jq_j=j\). Hence the two
constructions are inverse up to equivalence over \(P\).
\end{proof}

\subsection{Fixed points of fuzzy nuclei}
\label{subsec:fixed-points-fuzzy-nuclei}

\begin{lemma}[Intersections of fuzzy nuclei]
\label{lem:intersections-fuzzy-nuclei}
Let \((j_\alpha)_{\alpha\in I}\) be a family of fuzzy nuclei on \(P_\Lambda(X)\).
Then the pointwise meet
\[
    j(\phi)\eqdef\bigmeet_{\alpha\in I}j_\alpha(\phi)
\]
is again a fuzzy nucleus. If \(I=\varnothing\), this is the constant-top nucleus.
\end{lemma}

\begin{proof}
Inflationarity follows because \(\phi\leq j_\alpha(\phi)\) for every \(\alpha\).
Since every \(j_\alpha\) is monotone, \(j\) is monotone.

For idempotence, inflationarity gives
\[
    j(\phi)\leq j(j(\phi)).
\]
Conversely, since
\[
    j(\phi)\leq j_\alpha(\phi)
\]
for every \(\alpha\), monotonicity of \(j_\alpha\) gives
\[
    j_\alpha(j(\phi))
    \leq
    j_\alpha(j_\alpha(\phi))
    =
    j_\alpha(\phi).
\]
Taking the meet over \(\alpha\) yields
\[
    j(j(\phi))\leq j(\phi).
\]
Thus
\[
    j(j(\phi))=j(\phi).
\]

For finite meets,
\[
\begin{aligned}
j(\phi\meet\psi)
&=
\bigmeet_\alpha j_\alpha(\phi\meet\psi)
\\
&=
\bigmeet_\alpha\bigl(j_\alpha(\phi)\meet j_\alpha(\psi)\bigr)
\\
&=
\left(\bigmeet_\alpha j_\alpha(\phi)\right)
\meet
\left(\bigmeet_\alpha j_\alpha(\psi)\right)
\\
&=
j(\phi)\meet j(\psi).
\end{aligned}
\]

Finally, for scalar compatibility, inflationarity gives
\[
    r\tensor\phi\leq r\tensor j(\phi),
\]
and hence, by monotonicity of \(j\),
\[
    j(r\tensor\phi)\leq j(r\tensor j(\phi)).
\]
For the reverse inequality, since \(j(\phi)\leq j_\alpha(\phi)\), we have
\[
    r\tensor j(\phi)\leq r\tensor j_\alpha(\phi).
\]
Thus
\[
    j_\alpha(r\tensor j(\phi))
    \leq
    j_\alpha(r\tensor j_\alpha(\phi))
    =
    j_\alpha(r\tensor\phi).
\]
Taking the meet over \(\alpha\) gives
\[
    j(r\tensor j(\phi))
    \leq
    j(r\tensor\phi).
\]
Therefore
\[
    j(r\tensor j(\phi))=j(r\tensor\phi).
\]
\end{proof}

\begin{definition}[Fixed points of a fuzzy nucleus]
\label{def:fixed-points-fuzzy-nucleus}
Let \(j\) be a fuzzy nucleus on \(P_\Lambda(X)\). Define
\[
    P_\Lambda(X)_j
    \eqdef
    \{\phi\in P_\Lambda(X)\mid j(\phi)=\phi\}.
\]
For fixed points \(\phi_i,\phi,\psi\in P_\Lambda(X)_j\), define
\[
    \bigjoin\nolimits_j \phi_i
    =
    j\left(\bigjoin_i \phi_i\right),
\]
\[
    \phi\meet_j\psi
    =
    \phi\meet\psi,
\]
and
\[
    r\tensor_j\phi
    =
    j(r\tensor\phi).
\]
\end{definition}

\begin{proposition}[Fixed points of a fuzzy nucleus]
\label{prop:fixed-points-fuzzy-nucleus-frame}
Let \(j\) be a fuzzy nucleus on \(P_\Lambda(X)\). Then
\[
    P_\Lambda(X)_j
\]
is a frame. Its joins are closed joins,
\[
    \bigjoin\nolimits_j \phi_i
    =
    j\left(\bigjoin_i\phi_i\right),
\]
its finite meets are inherited from \(P_\Lambda(X)\), and it carries the closed scalar
action
\[
    r\tensor_j\phi
    =
    j(r\tensor\phi).
\]
The closed scalar action preserves arbitrary joins in each variable.
\end{proposition}

\begin{proof}
First, \(\top_{P_\Lambda(X)}\) is fixed by Lemma~\ref{lem:fuzzy-nuclei-preserve-top}.
If \(\phi,\psi\) are fixed, then
\[
    j(\phi\meet\psi)
    =
    j(\phi)\meet j(\psi)
    =
    \phi\meet\psi,
\]
so finite meets are inherited from \(P_\Lambda(X)\).

For a family \((\phi_i)\) of fixed points, define
\[
    c=j\left(\bigjoin_i\phi_i\right).
\]
Then \(c\) is fixed by idempotence. Since
\[
    \phi_i\leq \bigjoin_i\phi_i\leq j\left(\bigjoin_i\phi_i\right),
\]
the element \(c\) is an upper bound of the \(\phi_i\). If \(\psi\) is any fixed upper
bound of the \(\phi_i\), then
\[
    \bigjoin_i\phi_i\leq\psi.
\]
By monotonicity,
\[
    j\left(\bigjoin_i\phi_i\right)\leq j(\psi)=\psi.
\]
Thus \(c\) is the least fixed upper bound. Hence joins in \(P_\Lambda(X)_j\) are
closed joins. In particular, the bottom element is
\[
    \bot_j=j(\bot_{P_\Lambda(X)}).
\]

Frame distributivity follows from frame distributivity in \(P_\Lambda(X)\) and
finite-meet preservation by \(j\). If \(\phi\) and all \(\psi_i\) are fixed, then
\[
\begin{aligned}
\phi\meet_j\bigjoin\nolimits_j \psi_i
&=
\phi\meet j\left(\bigjoin_i \psi_i\right)
\\
&=
j(\phi)\meet j\left(\bigjoin_i \psi_i\right)
\\
&=
j\left(\phi\meet\bigjoin_i \psi_i\right)
\\
&=
j\left(\bigjoin_i(\phi\meet\psi_i)\right)
\\
&=
\bigjoin\nolimits_j(\phi\meet_j\psi_i).
\end{aligned}
\]

The closed scalar action is fixed by definition:
\[
    r\tensor_j\phi=j(r\tensor\phi).
\]
The scalar unit law follows from the unit law in \(P_\Lambda(X)\) and the fact that
\(\phi\) is fixed:
\[
    \top_L\tensor_j\phi
    =
    j(\top_L\tensor\phi)
    =
    j(\phi)
    =
    \phi.
\]
For associativity, use the scalar-compatibility equation:
\[
\begin{aligned}
r\tensor_j(s\tensor_j\phi)
&=
j(r\tensor j(s\tensor\phi))
\\
&=
j(r\tensor(s\tensor\phi))
\\
&=
j((r\tensor s)\tensor\phi)
\\
&=
(r\tensor s)\tensor_j\phi.
\end{aligned}
\]

It remains to verify that the closed scalar action preserves arbitrary joins. In the
frame variable,
\[
\begin{aligned}
r\tensor_j\bigjoin\nolimits_j \phi_i
&=
j\left(r\tensor j\left(\bigjoin_i\phi_i\right)\right)
\\
&=
j\left(r\tensor\bigjoin_i\phi_i\right)
\\
&=
j\left(\bigjoin_i(r\tensor\phi_i)\right).
\end{aligned}
\]
For any family \((\eta_i)\) in \(P_\Lambda(X)\), the closure identity
\[
    j\left(\bigjoin_i \eta_i\right)
    =
    j\left(\bigjoin_i j(\eta_i)\right)
\]
follows from inflationarity, monotonicity, and idempotence of \(j\). Applying this to
\(\eta_i=r\tensor\phi_i\), we get
\[
\begin{aligned}
j\left(\bigjoin_i(r\tensor\phi_i)\right)
&=
j\left(\bigjoin_i j(r\tensor\phi_i)\right)
\\
&=
\bigjoin\nolimits_j(r\tensor_j\phi_i).
\end{aligned}
\]
Therefore
\[
    r\tensor_j\bigjoin\nolimits_j \phi_i
    =
    \bigjoin\nolimits_j(r\tensor_j\phi_i).
\]

The scalar variable is similar. For a family \((r_i)\) in \(L\),
\[
\begin{aligned}
\left(\bigjoin_i r_i\right)\tensor_j\phi
&=
j\left(\left(\bigjoin_i r_i\right)\tensor\phi\right)
\\
&=
j\left(\bigjoin_i(r_i\tensor\phi)\right)
\\
&=
j\left(\bigjoin_i j(r_i\tensor\phi)\right)
\\
&=
\bigjoin\nolimits_j(r_i\tensor_j\phi).
\end{aligned}
\]
Thus the closed scalar action preserves arbitrary joins in each variable.
\end{proof}

\begin{remark}[Bottom scalars]
\label{rem:bottom-scalars}
Taking the empty join in the scalar variable gives
\[
    \bot_L\tensor_j\phi
    =
    \bot_j.
\]
Taking the empty join in the frame variable gives
\[
    r\tensor_j\bot_j
    =
    \bot_j.
\]
Thus each scalar operation preserves bottom, and the bottom scalar acts as the
constant-bottom operator. Notice that
\[
    \bot_j=j(\bot_{P_\Lambda(X)}),
\]
which need not equal the bottom element of \(P_\Lambda(X)\).
\end{remark}

\begin{lemma}[Enriched hom recovered from the scalar action]
\label{lem:enriched-hom-recovered-from-scalar-action}
Let
\[
    A=P_\Lambda(X)_j
\]
be the fixed-point frame of a fuzzy nucleus. For fixed points
\(\phi,\psi\in A\), one has
\[
    A(\phi,\psi)
    =
    \bigjoin
    \{r\in L\mid r\tensor_j\phi\leq \psi\}.
\]
Moreover this hom-value agrees with the hom inherited from \(P_\Lambda(X)\):
\[
    A(\phi,\psi)=P_\Lambda(X)(\phi,\psi).
\]
\end{lemma}

\begin{proof}
For \(r\in L\), using the tensor-hom formula in \(P_\Lambda(X)\), we have
\[
    r\leq P_\Lambda(X)(\phi,\psi)
    \quad\Longleftrightarrow\quad
    r\tensor\phi\leq \psi.
\]
Since \(\psi\) is fixed, this is equivalent to
\[
    j(r\tensor\phi)\leq \psi.
\]
But
\[
    j(r\tensor\phi)=r\tensor_j\phi.
\]
Hence
\[
    r\leq P_\Lambda(X)(\phi,\psi)
    \quad\Longleftrightarrow\quad
    r\tensor_j\phi\leq \psi.
\]
Taking the join over all such \(r\) gives the claim.
\end{proof}

\begin{definition}[\(\Lambda\)-fuzzy-frame presentation]
\label{def:fuzzy-frame-presentation}
A \textbf{\(\Lambda\)-fuzzy-frame presentation} is a pair
\[
    (X,j),
\]
where \(X\) is a fuzzy preorder and
\[
    j:P_\Lambda(X)\to P_\Lambda(X)
\]
is a fuzzy nucleus.

The frame presented by \((X,j)\) is the fixed-point frame
\[
    P_\Lambda(X)_j.
\]
\end{definition}

\begin{definition}[Fuzzy frame over \(\Lambda\)]
\label{def:fuzzy-frame}
A \textbf{fuzzy frame over \(\Lambda\)} is a frame presented by a
\(\Lambda\)-fuzzy-frame presentation. Equivalently, a fuzzy frame is a nuclear
quotient of a fuzzy presheaf frame:
\[
    P_\Lambda(X)\twoheadrightarrow P_\Lambda(X)_j.
\]
\end{definition}

Thus the phrase \emph{nuclear quotient} is frame-theoretic. Under locale duality, the
same structure is read as a fuzzy sublocale.

\begin{definition}[\(\Lambda\)-scalar frame]
\label{def:lambda-scalar-frame}
A \textbf{\(\Lambda\)-scalar frame} is a frame \(A\) equipped with an action
\[
    L\times A\to A,
    \qquad
    (r,a)\longmapsto r\tensor_A a,
\]
such that
\[
    \top_L\tensor_A a=a,
    \qquad
    (r\tensor s)\tensor_A a=r\tensor_A(s\tensor_A a),
\]
and such that the operation preserves arbitrary joins in each variable:
\[
    r\tensor_A\bigjoin_i a_i
    =
    \bigjoin_i(r\tensor_A a_i),
\]
\[
    \left(\bigjoin_i r_i\right)\tensor_A a
    =
    \bigjoin_i(r_i\tensor_A a).
\]
\end{definition}

Every fixed-point frame \(P_\Lambda(X)_j\) is a \(\Lambda\)-scalar frame by
Proposition~\ref{prop:fixed-points-fuzzy-nucleus-frame}.

\begin{definition}[Fuzzy-frame homomorphism]
\label{def:fuzzy-frame-homomorphism}
Let
\[
    P_\Lambda(X)_j
    \qquad\text{and}\qquad
    P_\Lambda(Y)_k
\]
be fuzzy frames presented by \(\Lambda\)-fuzzy-frame presentations. A
\textbf{fuzzy-frame homomorphism}
\[
    h:P_\Lambda(X)_j\to P_\Lambda(Y)_k
\]
is a frame homomorphism satisfying
\[
    h(r\tensor_j\phi)=r\tensor_k h(\phi)
\]
for all \(r\in L\) and all \(\phi\in P_\Lambda(X)_j\).
\end{definition}

\begin{corollary}[Fuzzy-frame homomorphisms are enriched functors]
\label{cor:fuzzy-frame-homomorphisms-are-enriched}
Let
\[
    h:A\to B
\]
be a fuzzy-frame homomorphism between fixed-point fuzzy frames. Then \(h\) is a
\(\Lambda\)-functor for the enriched homs recovered from the scalar actions.
\end{corollary}

\begin{proof}
Suppose
\[
    r\leq A(a,b).
\]
By Lemma~\ref{lem:enriched-hom-recovered-from-scalar-action}, this means
\[
    r\tensor_A a\leq b.
\]
Applying \(h\) and using scalar compatibility gives
\[
    r\tensor_B h(a)
    =
    h(r\tensor_A a)
    \leq
    h(b).
\]
Again by Lemma~\ref{lem:enriched-hom-recovered-from-scalar-action}, this implies
\[
    r\leq B(h(a),h(b)).
\]
Therefore
\[
    A(a,b)\leq B(h(a),h(b)),
\]
so \(h\) is a \(\Lambda\)-functor.
\end{proof}

\begin{definition}[The category of fuzzy frames]
\label{def:category-fuzzy-frames}
The category of fuzzy frames over \(\Lambda\) and fuzzy-frame homomorphisms is denoted
\[
    \FFrm_\Lambda.
\]
\end{definition}

\begin{remark}[Presentations]
\label{rem:fuzzy-frames-presentation-based}
Strictly speaking, one may keep the presentation \((X,j)\) as part of the data, or
identify presentations that present isomorphic fixed-point frames. In this chapter we
write fuzzy frames as the presented frames \(P_\Lambda(X)_j\), while remembering that
they arise from some localization of a fuzzy presheaf frame.

The definition is deliberately presentation-based. Fuzzy frames are not introduced as
arbitrary frames equipped with a scalar action; they are the fixed-point frames of
scalar-compatible nuclei on fuzzy presheaf frames.
\end{remark}

\begin{definition}[Nuclear quotient map]
\label{def:nuclear-quotient-map}
Let \(j\) be a fuzzy nucleus on \(P_\Lambda(X)\). The associated \textbf{nuclear
quotient map} is the closure map
\[
    q_j:
    P_\Lambda(X)\to P_\Lambda(X)_j,
    \qquad
    q_j(\phi)=j(\phi).
\]
\end{definition}

\begin{proposition}[The closure map is a quotient]
\label{prop:closure-map-is-quotient}
Let \(j\) be a fuzzy nucleus on \(P_\Lambda(X)\). Then
\[
    q_j:
    P_\Lambda(X)\to P_\Lambda(X)_j
\]
is a surjective fuzzy-frame homomorphism.
\end{proposition}

\begin{proof}
The map \(q_j\) is surjective because every fixed point \(\phi\in P_\Lambda(X)_j\)
satisfies
\[
    j(\phi)=\phi.
\]

It preserves arbitrary joins because joins in \(P_\Lambda(X)_j\) are closed joins:
\[
    j\left(\bigjoin_i\phi_i\right)
    =
    \bigjoin\nolimits_j j(\phi_i).
\]
For the empty family this says
\[
    q_j(\bot_{P_\Lambda(X)})=j(\bot_{P_\Lambda(X)})=\bot_j,
\]
so \(q_j\) also preserves bottom. It preserves finite meets because
\[
    j(\phi\meet\psi)=j(\phi)\meet j(\psi),
\]
and it preserves top by Lemma~\ref{lem:fuzzy-nuclei-preserve-top}.

It preserves scalars because
\[
\begin{aligned}
q_j(r\tensor\phi)
&=
j(r\tensor\phi)
\\
&=
j(r\tensor j(\phi))
\\
&=
r\tensor_j j(\phi)
\\
&=
r\tensor_j q_j(\phi).
\end{aligned}
\]
Therefore \(q_j\) is a surjective fuzzy-frame homomorphism.
\end{proof}

Conversely, every scalar-compatible quotient of a fuzzy presheaf frame determines a
fuzzy nucleus. This is the scalar-compatible version of the usual correspondence
between nuclei and quotient frames \cite{Johnstone1982}.

\begin{proposition}[Quotients determine fuzzy nuclei]
\label{prop:quotients-determine-fuzzy-nuclei}
Let \(B\) be a \(\Lambda\)-scalar frame, and let
\[
    q:P_\Lambda(X)\twoheadrightarrow B
\]
be a surjective frame homomorphism satisfying
\[
    q(r\tensor\phi)=r\tensor_B q(\phi)
\]
for all \(r\in L\) and \(\phi\in P_\Lambda(X)\). Since \(q\) preserves arbitrary
joins, it has a right adjoint
\[
    q_\ast:B\to P_\Lambda(X)
\]
defined by
\[
    q_\ast(b)
    =
    \bigjoin\{\phi\in P_\Lambda(X)\mid q(\phi)\leq b\}.
\]
Then
\[
    j_q\eqdef q_\ast q
\]
is a fuzzy nucleus on \(P_\Lambda(X)\), and the fixed-point fuzzy frame
\[
    P_\Lambda(X)_{j_q}
\]
is isomorphic to \(B\) as a fuzzy frame.
\end{proposition}

\begin{proof}
Since \(q\dashv q_\ast\), the composite
\[
    j_q=q_\ast q
\]
is inflationary. Since \(q\) is surjective, the counit
\[
    qq_\ast\leq \id_B
\]
is an equality. Indeed, if \(b=q(\phi)\), then \(\phi\leq q_\ast(b)\), so
\[
    b=q(\phi)\leq qq_\ast(b),
\]
while \(qq_\ast(b)\leq b\) follows from the counit inequality. Hence
\[
    qq_\ast=\id_B.
\]
Therefore
\[
    j_qj_q
    =
    q_\ast q q_\ast q
    =
    q_\ast q
    =
    j_q.
\]

Because \(q\) preserves finite meets and \(q_\ast\) preserves all meets,
\[
\begin{aligned}
j_q(\phi\meet\psi)
&=
q_\ast q(\phi\meet\psi)
\\
&=
q_\ast(q(\phi)\meet q(\psi))
\\
&=
q_\ast q(\phi)\meet q_\ast q(\psi)
\\
&=
j_q(\phi)\meet j_q(\psi).
\end{aligned}
\]

For scalar compatibility, use that \(q\) preserves scalars and that
\(qq_\ast=\id_B\):
\[
\begin{aligned}
j_q(r\tensor j_q(\phi))
&=
q_\ast q(r\tensor q_\ast q(\phi))
\\
&=
q_\ast\bigl(r\tensor_B q q_\ast q(\phi)\bigr)
\\
&=
q_\ast(r\tensor_B q(\phi))
\\
&=
q_\ast q(r\tensor\phi)
\\
&=
j_q(r\tensor\phi).
\end{aligned}
\]
Therefore \(j_q\) is a fuzzy nucleus.

The restriction
\[
    q:P_\Lambda(X)_{j_q}\to B
\]
is an isomorphism of frames with inverse \(q_\ast\). First, for every \(b\in B\),
\[
    j_q(q_\ast b)
    =
    q_\ast q q_\ast b
    =
    q_\ast b,
\]
so \(q_\ast b\) is fixed. Moreover,
\[
    q q_\ast(b)=b,
\]
and for \(\phi\in P_\Lambda(X)_{j_q}\),
\[
    q_\ast q(\phi)=j_q(\phi)=\phi.
\]

It remains to note that this isomorphism respects the scalar actions. The restriction
of \(q\) preserves scalars by assumption. Conversely, for \(b\in B\), write
\(b=q(\phi)\). Then
\[
\begin{aligned}
q\bigl(r\tensor_{j_q}q_\ast(b)\bigr)
&=
q\bigl(j_q(r\tensor q_\ast(b))\bigr)
\\
&=
q(r\tensor q_\ast(b))
\\
&=
r\tensor_B qq_\ast(b)
\\
&=
r\tensor_B b.
\end{aligned}
\]
Since \(q:P_\Lambda(X)_{j_q}\to B\) is an isomorphism, this implies
\[
    r\tensor_{j_q}q_\ast(b)
    =
    q_\ast(r\tensor_B b).
\]
Thus
\[
    P_\Lambda(X)_{j_q}\cong B
\]
as fuzzy frames.
\end{proof}

\subsection{Examples}
\label{subsec:localizing-presheaves-examples}

\begin{example}[The identity localization]
\label{ex:identity-localization}
For every fuzzy preorder \(X\), the identity map
\[
    \id_{P_\Lambda(X)}:
    P_\Lambda(X)\to P_\Lambda(X)
\]
is a fuzzy nucleus. Hence
\[
    P_\Lambda(X)_{\id}
    =
    P_\Lambda(X).
\]
Thus every fuzzy presheaf frame is itself a fuzzy frame.

In this case no predicates are identified or closed further. Joins, finite meets, and
scalar multiplication are the pointwise operations inherited directly from
\(\Lambda\). Localically, this is the ambient presheaf fuzzy locale
\[
    \mathcal P_\Lambda X.
\]
All other examples of fuzzy frames in this chapter may be read as sublocales of such
presheaf fuzzy locales, equivalently as nuclear quotients of their frames of opens.
\end{example}

\begin{example}[The trivial localization]
\label{ex:terminal-localization}
For every fuzzy preorder \(X\), the constant-top map
\[
    j_\top(\phi)=\top_{P_\Lambda(X)}
\]
is a fuzzy nucleus. Inflationarity and idempotence are immediate, and
\[
    j_\top(\phi\meet\psi)=\top_{P_\Lambda(X)}
    =
    \top_{P_\Lambda(X)}\meet\top_{P_\Lambda(X)}.
\]
For scalar compatibility,
\[
    j_\top(r\tensor j_\top(\phi))
    =
    \top_{P_\Lambda(X)}
    =
    j_\top(r\tensor\phi).
\]
The fixed-point frame has exactly one element:
\[
    P_\Lambda(X)_{j_\top}\cong 1.
\]
This is the trivial fuzzy frame. On the locale side it corresponds to the initial, or
empty, fuzzy sublocale of the presheaf fuzzy locale. The identity nucleus gives the
whole sublocale, while the constant-top nucleus gives the least sublocale.
\end{example}

\begin{example}[The crisp case]
\label{ex:crisp-localization-case}
Let
\[
    \Lambda=2.
\]
Then fuzzy preorders are ordinary preorders and, for an ordinary preorder \(X\),
\[
    P_2(X)
    =
    [X^{\op},2]
    \cong
    \Down(X),
\]
the frame of downsets of \(X\). A fuzzy nucleus on \(P_2(X)\) is exactly an ordinary
nucleus on the frame \(\Down(X)\), because the scalar action of \(2\) is forced:
\[
    \top\tensor U=U,
    \qquad
    \bot\tensor U=\bot.
\]
Thus the scalar-compatibility equation adds no extra condition beyond the ordinary
nucleus axioms.

Concretely, a nucleus
\[
    j:\Down(X)\to\Down(X)
\]
assigns to each downset \(U\subseteq X\) a closed downset \(j(U)\), satisfying
\[
    U\subseteq j(U),
    \qquad
    j(j(U))=j(U),
    \qquad
    j(U\cap V)=j(U)\cap j(V).
\]
The fixed points form the ordinary quotient frame
\[
    \Down(X)_j.
\]
Thus the crisp specialization recovers the standard description of frames as
nuclear quotients of downset frames.
\end{example}

\begin{theorem}[Ordinary frames as localized presheaves]
\label{thm:ordinary-frames-localized-presheaves}
Let \(A\) be a small ordinary frame. Regard the underlying poset of \(A\) as a
preorder. Then the map
\[
    q:\Down(A)\to A,
    \qquad
    q(U)=\bigjoin U,
\]
is a surjective frame homomorphism. Consequently \(A\) is isomorphic to the
fixed-point frame of the nucleus
\[
    j_A=q_\ast q
\]
on \(\Down(A)\). Explicitly,
\[
    j_A(U)=\downarrow\bigjoin U.
\]
\end{theorem}

\begin{proof}
The map \(q\) preserves arbitrary joins because
\[
\begin{aligned}
q\left(\bigcup_i U_i\right)
&=
\bigjoin\bigcup_i U_i
\\
&=
\bigjoin_i\bigjoin U_i
\\
&=
\bigjoin_i q(U_i).
\end{aligned}
\]
It preserves top because
\[
    q(A)=\bigjoin A=\top_A.
\]

It remains to check finite meets. Since
\[
    U\cap V\subseteq U
    \qquad\text{and}\qquad
    U\cap V\subseteq V,
\]
we have
\[
    q(U\cap V)\leq q(U)\meet q(V).
\]
Conversely, by frame distributivity,
\[
\begin{aligned}
q(U)\meet q(V)
&=
\left(\bigjoin U\right)\meet\left(\bigjoin V\right)
\\
&=
\bigjoin_{u\in U,\;v\in V}(u\meet v).
\end{aligned}
\]
Since \(U\) and \(V\) are downsets, for every \(u\in U\) and \(v\in V\), the element
\[
    u\meet v
\]
belongs to \(U\cap V\). Hence
\[
    q(U)\meet q(V)\leq q(U\cap V).
\]
Therefore
\[
    q(U\cap V)=q(U)\meet q(V).
\]

Finally, \(q\) is surjective because
\[
    q(\downarrow a)=a
\]
for every \(a\in A\). Therefore \(q\) is a surjective frame homomorphism. Its right
adjoint is
\[
    q_\ast(a)=\downarrow a.
\]
Hence
\[
    j_A(U)
    =
    q_\ast q(U)
    =
    \downarrow\bigjoin U.
\]
Proposition~\ref{prop:quotients-determine-fuzzy-nuclei} gives
\[
    \Down(A)_{j_A}\cong A.
\]
\end{proof}

\begin{theorem}[The crisp case recovers ordinary frames]
\label{thm:crisp-case-recovers-ordinary-frames}
When
\[
    \Lambda=2,
\]
the small fuzzy frames over \(2\) are exactly the small ordinary frames, up to
isomorphism.
\end{theorem}

\begin{proof}
Every fuzzy frame over \(2\) is, by definition, a nuclear quotient of a frame
\[
    \Down(X).
\]
By Example~\ref{ex:crisp-localization-case}, this is an ordinary nuclear quotient of
an ordinary frame, and hence is an ordinary frame.

Conversely, every small ordinary frame is a nuclear quotient of an ordinary presheaf
frame by Theorem~\ref{thm:ordinary-frames-localized-presheaves}. Therefore every
small ordinary frame is a fuzzy frame over \(2\).
\end{proof}

This recovers the usual pointfree frame-theoretic localization picture
\cite{Johnstone1982}.

\begin{example}[G\"odel-valued localizations]
\label{ex:godel-valued-localizations}
Let
\[
    \Lambda=[0,1]
\]
with G\"odel tensor
\[
    r\tensor s=r\meet s=\min(r,s).
\]
A fuzzy predicate on a fuzzy preorder \(X\) is a map
\[
    \phi:X\to[0,1]
\]
satisfying
\[
    X(x',x)\meet\phi(x)\leq \phi(x')
\]
for all \(x,x'\in X\). The scalar action is truncation by the level \(r\):
\[
    (r\tensor\phi)(x)=\min(r,\phi(x)).
\]
A G\"odel-valued fuzzy nucleus closes predicates compatibly with truncation:
\[
    j(r\meet j(\phi))=j(r\meet\phi).
\]

Thus if a predicate has already been closed, then truncating it at level \(r\) and
closing again gives the same result as truncating the original predicate and closing.
The fixed-point frame consists of the predicates satisfying the chosen closure
condition. Joins are closed pointwise suprema, finite meets are pointwise minima, and
the scalar action is closed truncation:
\[
    r\tensor_j\phi=j(r\meet\phi).
\]
This gives a localic form of G\"odel-valued fuzzy topology, obtained here by enriched
presheaf localization rather than by postulating a topology of fuzzy opens directly
\cite{HohleRodabaugh1999}.
\end{example}

\begin{example}[{\L}ukasiewicz-valued localizations]
\label{ex:lukasiewicz-valued-localizations}
Let
\[
    \Lambda=[0,1]
\]
with {\L}ukasiewicz tensor
\[
    r\tensor s=\max(0,r+s-1).
\]
A fuzzy predicate on \(X\) is a map
\[
    \phi:X\to[0,1]
\]
satisfying
\[
    X(x',x)\tensor\phi(x)\leq \phi(x').
\]
The scalar action is no longer truncation. It shifts truth degrees downward:
\[
    (r\tensor\phi)(x)=\max(0,r+\phi(x)-1).
\]
For \(r=1\), this is the identity; for \(r<1\), all values are lowered, and values
below \(1-r\) are sent to \(0\).

A {\L}ukasiewicz-valued fuzzy nucleus is therefore a closure operator satisfying
\[
    j\bigl(r\tensor j(\phi)\bigr)=j(r\tensor\phi).
\]
The fixed-point fuzzy frame has closed joins, inherited finite meets, and closed
{\L}ukasiewicz scalar action. This example emphasizes that the scalar structure is not
merely order-theoretic: the same underlying lattice \([0,1]\) produces different
fuzzy-localic behavior depending on the chosen tensor.
\end{example}

\begin{example}[Product-valued localizations]
\label{ex:product-valued-localizations}
Let
\[
    \Lambda=[0,1]
\]
with product tensor
\[
    r\tensor s=rs.
\]
Then a fuzzy predicate on \(X\) is a map
\[
    \phi:X\to[0,1]
\]
satisfying
\[
    X(x',x)\phi(x)\leq \phi(x')
\]
for all \(x,x'\in X\). Scalar multiplication is ordinary multiplication:
\[
    (r\tensor\phi)(x)=r\phi(x).
\]
A product-valued fuzzy nucleus closes these predicates compatibly with product
rescaling:
\[
    j(r\,j(\phi))=j(r\,\phi).
\]

In the fixed-point frame, multiplying a closed predicate by \(r\) may fail to remain
closed, so the closed scalar action is
\[
    r\tensor_j\phi=j(r\phi).
\]
This can be interpreted as first weakening the degree of membership by the factor
\(r\), and then imposing the local closure condition. Product-valued localizations
therefore behave differently from G\"odel-valued localizations: G\"odel scalars cut
truth values from above, whereas product scalars contract them proportionally.
\end{example}

\begin{example}[Lawvere-valued localizations]
\label{ex:lawvere-valued-localizations}
Let
\[
    \Lambda=
    ([0,\infty],\leq_L,\meet_L,\join_L,+,\ominus,\infty,0),
    \qquad
    a\leq_L b
    \Longleftrightarrow
    a\geq b
    \text{ in the usual order}.
\]
Thus the tensor is addition, the top element is \(0\), and the bottom element is
\(\infty\). Joins in \(\Lambda\) are infima in the usual order, and meets in
\(\Lambda\) are suprema in the usual order.

A fuzzy preorder is a Lawvere metric space. A Lawvere-valued fuzzy predicate is a map
\[
    \phi:X\to[0,\infty]
\]
satisfying, in the usual order,
\[
    \phi(x')\leq d(x',x)+\phi(x).
\]
Thus \(\phi(x)\) is better read as a cost, distance, or potential rather than as a
truth degree. The scalar action adds a constant:
\[
    (r\tensor\phi)(x)=r+\phi(x).
\]
Since the order is reversed, inflationarity of a nucleus
\[
    \phi\leq_L j(\phi)
\]
means
\[
    j(\phi)(x)\leq \phi(x)
\]
in the usual order. Thus a Lawvere-valued nucleus lowers costs, subject to the chosen
closure condition.

The scalar-compatibility equation becomes
\[
    j(r+j(\phi))=j(r+\phi),
\]
where the displayed equation is written in the usual arithmetic notation on
\([0,\infty]\). The fixed-point frame consists of the closed cost predicates. Closed
joins are closures of pointwise infima, while finite meets are pointwise suprema.
This is the metric analogue of passing from arbitrary presheaves to local objects
\cite{Lawvere1973MetricSpaces}.
\end{example}

\section{Fuzzy locales}
\label{sec:fuzzy-locales}

We now pass from fuzzy frames to fuzzy locales by formal duality. As in ordinary
pointfree topology, frame homomorphisms are read contravariantly as inverse-image
maps of locales \cite{Johnstone1982}.

\subsection{Fuzzy locales as duals of fuzzy frames}
\label{subsec:fuzzy-locales-duals}

\begin{definition}[Fuzzy locale]
\label{def:fuzzy-locale}
The category of \textbf{fuzzy locales over \(\Lambda\)} is
\[
    \FLoc_\Lambda
    \eqdef
    \FFrm_\Lambda^{\op}.
\]
If \(\mathcal X\) is a fuzzy locale, its fuzzy frame of opens is denoted
\[
    \Opens_\Lambda(\mathcal X).
\]
A map of fuzzy locales
\[
    f:\mathcal X\to\mathcal Y
\]
is represented by a fuzzy-frame homomorphism in the opposite direction:
\[
    f^\ast:
    \Opens_\Lambda(\mathcal Y)
    \to
    \Opens_\Lambda(\mathcal X).
\]
\end{definition}

Thus the fundamental duality is
\[
\begin{tikzcd}[column sep=large, row sep=large]
\mathcal X
    \arrow[r, "f"]
&
\mathcal Y
\\
\Opens_\Lambda(\mathcal X)
&
\Opens_\Lambda(\mathcal Y)
    \arrow[l, "f^\ast"']
\end{tikzcd}
\]
and, categorically,
\[
\begin{tikzcd}[column sep=large]
\FLoc_\Lambda
    \arrow[r, "\Opens_\Lambda"]
&
\FFrm_\Lambda^{\op}.
\end{tikzcd}
\]

\subsection{Locales presented by fuzzy presheaves}
\label{subsec:locales-presented-by-fuzzy-presheaves}

For a fuzzy preorder \(X\), the fuzzy presheaf frame \(P_\Lambda(X)\) determines a
fuzzy locale.

\begin{definition}[Presheaf fuzzy locale]
\label{def:presheaf-fuzzy-locale}
For a fuzzy preorder \(X\), the \textbf{presheaf fuzzy locale} associated to \(X\) is
the fuzzy locale
\[
    \mathcal P_\Lambda X
\]
whose fuzzy frame of opens is
\[
    \Opens_\Lambda(\mathcal P_\Lambda X)
    =
    P_\Lambda(X).
\]
Equivalently, \(\mathcal P_\Lambda X\) is the fuzzy locale corresponding to the
identity nucleus on \(P_\Lambda(X)\).
\end{definition}

\begin{proposition}[Functoriality of presheaf fuzzy locales]
\label{prop:functoriality-presheaf-fuzzy-locales}
A \(\Lambda\)-functor
\[
    f:X\to Y
\]
induces a map of fuzzy locales
\[
    \mathcal P_\Lambda X\to\mathcal P_\Lambda Y.
\]
\end{proposition}

\begin{proof}
Precomposition gives
\[
    f^\ast:P_\Lambda(Y)\to P_\Lambda(X),
    \qquad
    f^\ast(\psi)=\psi f.
\]
Explicitly, if \(\psi:Y\to L\) is a fuzzy predicate on \(Y\), then
\[
\begin{aligned}
X(x',x)\tensor(\psi f)(x)
&=
X(x',x)\tensor\psi(fx)
\\
&\leq
Y(fx',fx)\tensor\psi(fx)
\\
&\leq
\psi(fx')
\\
&=
(\psi f)(x').
\end{aligned}
\]
Thus \(f^\ast(\psi)\in P_\Lambda(X)\). Since precomposition preserves pointwise
joins, finite meets, and scalars, \(f^\ast\) is a fuzzy-frame homomorphism. Therefore
it defines a fuzzy-locale map
\[
    \mathcal P_\Lambda X\to\mathcal P_\Lambda Y.
\]
\end{proof}

The construction is covariant on fuzzy preorders:
\[
\begin{tikzcd}[column sep=large, row sep=large]
X
    \arrow[r, "f"]
    \arrow[d, "\mathcal P_\Lambda"']
&
Y
    \arrow[d, "\mathcal P_\Lambda"]
\\
\mathcal P_\Lambda X
    \arrow[r]
&
\mathcal P_\Lambda Y
\end{tikzcd}
\]
while the inverse-image map on opens is contravariant:
\[
\begin{tikzcd}[column sep=large, row sep=large]
\Opens_\Lambda(\mathcal P_\Lambda X)
&
\Opens_\Lambda(\mathcal P_\Lambda Y)
    \arrow[l, "f^\ast"']
\\
P_\Lambda(X)
    \arrow[u, equal]
&
P_\Lambda(Y)
    \arrow[l, "f^\ast"]
    \arrow[u, equal]
\end{tikzcd}
\]

\subsection{Localized presheaves as fuzzy sublocales}
\label{subsec:localized-presheaves-fuzzy-sublocales}

A fuzzy nucleus has two equivalent readings. Frame-theoretically, it is a quotient
of the fuzzy presheaf frame:
\[
    P_\Lambda(X)\twoheadrightarrow P_\Lambda(X)_j.
\]
Localically, it is a sublocale of the presheaf fuzzy locale. This is the ordinary
duality between nuclei, quotient frames, and sublocales, with the scalar compatibility
carried along \cite{Johnstone1982}.

\begin{definition}[Fuzzy sublocale presented by a nucleus]
\label{def:fuzzy-sublocale-presented-by-nucleus}
Let \(j\) be a fuzzy nucleus on \(P_\Lambda(X)\). The \textbf{fuzzy sublocale
presented by \(j\)} is the fuzzy locale
\[
    \mathcal X_j
\]
whose frame of opens is
\[
    \Opens_\Lambda(\mathcal X_j)
    =
    P_\Lambda(X)_j.
\]
\end{definition}

The quotient and sublocale readings are displayed by
\[
\begin{tikzcd}[column sep=large, row sep=large]
P_\Lambda(X)
    \arrow[r, two heads, "q_j"]
&
P_\Lambda(X)_j
\\
\mathcal P_\Lambda X
&
\mathcal X_j
    \arrow[l, hook]
\end{tikzcd}
\]
where the top row is in fuzzy frames and the bottom row is in fuzzy locales. The hook
on the locale side is justified by the fact that \(q_j\) is a quotient frame
homomorphism.

\subsection{Examples}
\label{subsec:fuzzy-locale-examples}

\begin{example}[Ordinary presheaf locales]
\label{ex:ordinary-presheaf-locales}
Let
\[
    \Lambda=2.
\]
For an ordinary preorder \(X\),
\[
    \Opens_2(\mathcal P_2X)
    =
    P_2(X)
    =
    \Down(X).
\]
Thus \(\mathcal P_2X\) is the ordinary lower presheaf locale associated to \(X\).
Nuclei on \(\Down(X)\) correspond to ordinary sublocales of this locale
\cite{Johnstone1982}.

If \(X\) is discrete, then \(\Down(X)\cong \mathcal P(X)\), so the corresponding
locale is the ordinary powerset locale on the set \(X\). If \(X\) has nontrivial
order, then the opens are downward-closed subsets, and the locale remembers the
Alexandrov topology determined by the preorder.
\end{example}

\begin{example}[G\"odel-valued presheaf locales]
\label{ex:godel-valued-presheaf-locales}
For G\"odel truth values,
\[
    \Opens_\Lambda(\mathcal P_\Lambda X)
    =
    P_\Lambda(X)
\]
is the frame of G\"odel-valued fuzzy predicates on \(X\). A fuzzy open is therefore a
monotone fuzzy lower set, where monotonicity is measured by the G\"odel tensor:
\[
    X(x',x)\meet\phi(x)\leq \phi(x').
\]
A fuzzy nucleus on this frame presents a G\"odel-valued fuzzy sublocale. Its opens are
exactly the closed G\"odel predicates. Localically, the quotient
\[
    P_\Lambda(X)\twoheadrightarrow P_\Lambda(X)_j
\]
is read as the inclusion of the fuzzy sublocale
\[
    \mathcal X_j\hookrightarrow \mathcal P_\Lambda X.
\]
\end{example}

\begin{example}[{\L}ukasiewicz-valued presheaf locales]
\label{ex:lukasiewicz-valued-presheaf-locales}
For {\L}ukasiewicz truth values,
\[
    \Opens_\Lambda(\mathcal P_\Lambda X)
    =
    P_\Lambda(X)
\]
is the frame of {\L}ukasiewicz-valued fuzzy predicates on \(X\). The enriched
lower-set condition is
\[
    \max(0,X(x',x)+\phi(x)-1)\leq \phi(x').
\]
A fuzzy nucleus presents a localized fuzzy locale whose opens are the corresponding
closed predicates. The scalar action lowers degrees by {\L}ukasiewicz multiplication
and then recloses them:
\[
    r\tensor_j\phi
    =
    j\bigl(\max(0,r+\phi-1)\bigr),
\]
where the formula is interpreted pointwise.
\end{example}

\begin{example}[Product-valued presheaf locales]
\label{ex:product-valued-presheaf-locales}
For product truth values,
\[
    \Opens_\Lambda(\mathcal P_\Lambda X)
    =
    P_\Lambda(X)
\]
is the frame of product-valued fuzzy predicates satisfying
\[
    X(x',x)\phi(x)\leq \phi(x').
\]
A fuzzy nucleus gives a product-valued fuzzy sublocale whose opens are closed under
the chosen closure operation. The scalar action multiplies fuzzy predicates by
constants and then closes:
\[
    r\tensor_j\phi=j(r\phi).
\]
Thus scalar rescaling behaves like proportional attenuation of fuzzy membership.
\end{example}

\begin{example}[Lawvere-valued presheaf locales]
\label{ex:lawvere-valued-presheaf-locales}
For a Lawvere metric space \(X\), the presheaf fuzzy locale \(\mathcal P_\Lambda X\)
has frame of opens
\[
    \Opens_\Lambda(\mathcal P_\Lambda X)=P_\Lambda(X),
\]
whose elements are maps
\[
    \phi:X\to[0,\infty]
\]
satisfying, in the usual order,
\[
    \phi(x')\leq d(x',x)+\phi(x).
\]
These are nonexpansive cost predicates. A nucleus imposes a local closure condition
on such cost predicates. Since the order is reversed, closed joins correspond to
closures of pointwise infima of costs, while finite meets correspond to pointwise
suprema. This gives a metric-localic interpretation of Lawvere presheaves.
\end{example}

\section{The fuzzy Sierpiński classifier}
\label{sec:fuzzy-sierpinski-classifier}

In ordinary locale theory, the Sierpiński locale classifies opens. On the frame
side, this says that the Sierpiński frame is freely generated by one open. In the
present enriched setting the same principle holds, but the free object is not, in
general, the raw presheaf category on the two-object subterminal shape. For
non-idempotent fuzzy tensors, such as the {\L}ukasiewicz and product tensors, that
raw two-point construction does not classify arbitrary fuzzy opens by
finite-limit-preserving fuzzy-frame morphisms.

The correct construction uses a small enriched syntactic sketch of finite fuzzy-open
formulae and then takes the associated left-exact enriched-sketch localization of its
fuzzy presheaf category. The resulting object is characterized by its classifying
universal property.

\subsection{Categorical fuzzy frames and fuzzy opens}
\label{subsec:categorical-fuzzy-frames-and-fuzzy-opens}

We recall the categorical form of fuzzy frames used in this chapter. A categorical
fuzzy frame is a left-exact reflective localization
\[
    a\dashv i:
    P_{\mathbb V}(X)
    \rightleftarrows
    \mathcal A
\]
of a fuzzy presheaf category, where \(X\) is a small \(\mathbb V\)-category,
\(i\) is fully faithful, and \(a\) preserves the chosen finite
\(\mathcal F\)-weighted limits.

A morphism of categorical fuzzy frames
\[
    H:\mathcal A\to\mathcal B
\]
is a \(\mathbb V\)-functor preserving small weighted colimits and the chosen finite
\(\mathcal F\)-weighted limits. Enriched natural transformations are the 2-cells.
We write
\[
    \mathbf{FFrm}_{\mathbb V}
\]
for the resulting 2-category: objects are categorical fuzzy frames, 1-cells are
small-colimit-preserving and finite-limit-preserving \(\mathbb V\)-functors, and
2-cells are enriched natural transformations. Thus
\[
    \mathbf{FFrm}_{\mathbb V}(\mathcal A,\mathcal B)
\]
denotes the ordinary hom-category of such morphisms and enriched natural
transformations.

A cosmos-level fuzzy open of a categorical fuzzy frame \(\mathcal A\) is a subobject
\[
    S\hookrightarrow 1_{\mathcal A}
\]
of the terminal object in the underlying ordinary category of \(\mathcal A\). We
write
\[
    \operatorname{Sub}_{\mathcal A}(1_{\mathcal A})
\]
for the ordinary category of such subobjects. Thus, if \(\mathcal A\) is enriched,
subobjects are computed in the underlying ordinary category \(\mathcal A_0\).

Equivalently, a representative object \(S\) is subterminal precisely when the
diagonal
\[
    S\longrightarrow S\times S
\]
is an isomorphism. Since morphisms of categorical fuzzy frames preserve terminal
objects and binary products, they send subterminal objects to subterminal objects.

We shall use the following tensor calculation for categorical fuzzy frames.

\begin{lemma}[Tensors in categorical fuzzy frames]
\label{lem:categorical-fuzzy-frames-are-tensored}
Let
\[
    a\dashv i:
    P_{\mathbb V}(X)\rightleftarrows\mathcal A
\]
be a categorical fuzzy-frame presentation. Then \(\mathcal A\) is tensored over
\(\mathbb V\). For \(V\in\mathbb V\) and \(A\in\mathcal A\), the tensor is computed
by
\[
    V\tensor_{\mathcal A}A
    \cong
    a(V\tensor iA).
\]
\end{lemma}

\begin{proof}
For \(B\in\mathcal A\), using the enriched adjunction and the tensor--hom
adjunction in \(P_{\mathbb V}(X)\), we compute
\[
\begin{aligned}
\mathcal A(a(V\tensor iA),B)
&\cong
P_{\mathbb V}(X)(V\tensor iA,iB)
\\
&\cong
[V,P_{\mathbb V}(X)(iA,iB)]
\\
&\cong
[V,\mathcal A(A,B)].
\end{aligned}
\]
Thus \(a(V\tensor iA)\) represents the tensor of \(A\) by \(V\) in
\(\mathcal A\).
\end{proof}

\subsection{The enriched syntactic sketch of one fuzzy open}
\label{subsec:enriched-syntactic-sketch-one-fuzzy-open}

We now construct the enriched syntactic sketch whose models are subterminal objects
together with the formal finite-limit and scalar formulae generated from them.

The construction is made relative to the universe convention fixed earlier. In
particular, the collection of scalar objects used in the sketch, the chosen class
\(\mathcal F\) of finite weights, and the generated syntax below are small in the
ambient universe in which enriched sketches are formed. When the fuzzy cosmos
\(\mathbb V\) is not small in the lowest universe, this construction is understood
in the next universe of the fixed tower, equivalently after choosing a small
skeleton of the scalar objects used by the sketch.

Let
\[
    \mathbb T_{\Omega}
\]
denote the small enriched finite-limit-and-scalar sketch freely generated by one
subterminal object. More explicitly, \(\mathbb T_{\Omega}\) is generated as follows.

First, there are two basic formal objects
\[
    1,
    \qquad
    U.
\]
The object \(1\) is declared terminal, and \(U\) is declared subterminal. Thus the
sketch contains the terminal-object data for \(1\), the canonical map
\[
    U\longrightarrow 1,
\]
and the finite-limit data expressing that the diagonal comparison
\[
    U\longrightarrow U\times U
\]
is an isomorphism.

Second, the sketch is closed under the chosen finite \(\mathcal F\)-weighted limit
operations. Thus, whenever a formal finite-limit expression
\[
    \{\varphi,D\}
\]
is built from earlier formal objects, it is again a formal object of
\(\mathbb T_{\Omega}\), together with its designated limiting cone. In particular,
binary products of formal objects are again formal objects.

Third, for every scalar object \(V\in\mathbb V\) and every formal object
\(A\in\mathbb T_{\Omega}\), there is a formal scalar object
\[
    V\diamond A.
\]
The symbol \(V\diamond A\) is syntactic at this stage. It is not required to
represent the tensor of \(A\) by \(V\) in the raw presheaf category, nor is it
required to satisfy a tensor universal property inside \(\mathbb T_{\Omega}\).

For each formal scalar object \(V\diamond A\), the sketch contains a generating
enriched morphism
\[
    \kappa_{V,A}:
    V\longrightarrow
    \mathbb T_{\Omega}(A,V\diamond A).
\]
Equivalently, \(\kappa_{V,A}\) is a formal tensor coevaluation map. Under Yoneda,
it determines a canonical comparison map in the presheaf category
\[
    \alpha_{V,A}:
    V\tensor yA
    \longrightarrow
    y(V\diamond A),
\]
where
\[
    y:\mathbb T_{\Omega}\to
    P_{\mathbb V}(\mathbb T_{\Omega})
\]
is the enriched Yoneda embedding.

The hom-objects and equations of \(\mathbb T_{\Omega}\) are generated freely by the
terminal-object data, the subterminality data for \(U\), the specified finite-limit
cones, the scalar-forming operations \(V\diamond A\), and the generating maps
\(\kappa_{V,A}\), subject only to the coherence equations for the specified
finite-limit cones and for functoriality of the displayed generating data. No tensor
universal property for \(V\diamond A\) is imposed inside
\(\mathbb T_{\Omega}\). The tensor universal property is imposed only on models, by
requiring the transposes
\[
    V\tensor_{\mathcal A}F(A)\longrightarrow F(V\diamond A)
\]
to be isomorphisms.

Equivalently, \(\mathbb T_{\Omega}\) is the free small enriched
finite-limit-and-scalar sketch generated by one subterminal object \(U\), where the
scalar operations are formal operations whose intended interpretation is by tensors
in the target category of models.

\begin{remark}[Formal scalar objects]
\label{rem:formal-scalar-objects}
One should not impose
\[
    y(V\diamond A)\cong V\tensor yA
\]
inside the raw presheaf category. In general, a scalar multiple of a representable
presheaf need not be representable.

In the posetal specialization, if such an identity held representably, then for a
scalar \(v<\top_L\) it would force
\[
    \mathbb T_{\Omega}(Z,V\diamond A)
    =
    v\tensor \mathbb T_{\Omega}(Z,A).
\]
Taking \(Z=V\diamond A\) would give
\[
    \top_L
    \leq
    \mathbb T_{\Omega}(V\diamond A,V\diamond A)
    =
    v\tensor \mathbb T_{\Omega}(V\diamond A,A)
    \leq
    v,
\]
which is impossible. Thus \(V\diamond A\) is a formal scalar formula, and the
comparison
\[
    V\tensor yA\to y(V\diamond A)
\]
is imposed only after passing to the sketch localization, and only in the correct
classifying variance.
\end{remark}

\subsection{Models of the Sierpiński sketch}
\label{subsec:models-sierpinski-sketch}

Let \(\mathcal A\) be a categorical fuzzy frame. By
Lemma~\ref{lem:categorical-fuzzy-frames-are-tensored}, \(\mathcal A\) is tensored
over \(\mathbb V\). If
\[
    a\dashv i:
    P_{\mathbb V}(X)\rightleftarrows\mathcal A
\]
is a presentation of \(\mathcal A\), then the tensor of \(A\in\mathcal A\) by
\(V\in\mathbb V\) is computed by
\[
    V\tensor_{\mathcal A}A
    \cong
    a(V\tensor iA).
\]

A \textbf{model of the Sierpiński sketch} in \(\mathcal A\) is a
\(\mathbb V\)-functor
\[
    F:\mathbb T_{\Omega}\to\mathcal A
\]
which satisfies the following conditions:
\begin{enumerate}
    \item \(F\) sends the specified finite \(\mathcal F\)-limit cones to finite
    limits in \(\mathcal A\);

    \item \(F(U)\) is subterminal;

    \item for every scalar object \(V\in\mathbb V\) and every formal object
    \(A\in\mathbb T_{\Omega}\), the image of the generating morphism
    \[
        \kappa_{V,A}:
        V\to \mathbb T_{\Omega}(A,V\diamond A)
    \]
    has transpose
    \[
        \overline{\kappa}^{F}_{V,A}:
        V\tensor_{\mathcal A}F(A)\to F(V\diamond A),
    \]
    and this transpose is an isomorphism.
\end{enumerate}
Thus scalar formulae are interpreted as actual tensors in \(\mathcal A\), and the
formal coevaluation maps are interpreted as the corresponding tensor
coevaluations.

\begin{remark}[Scalar formulae need not be subterminal]
\label{rem:scalar-formulae-not-necessarily-subterminal}
Although the generator \(U\) is required to be subterminal, the formal scalar objects
\(V\diamond A\) are not declared subterminal. In a general categorical fuzzy frame,
tensors of subterminal objects need not be subterminal. Thus the cosmos-level
classifier remembers one distinguished fuzzy open
\[
    U\hookrightarrow 1,
\]
while the scalar and finite-limit formulae are formal objects generated from it. In
the posetal reflection every object of a frame lies below the top element, so these
formula objects become ordinary fuzzy opens.
\end{remark}

We write
\[
    \operatorname{Mod}_{\Omega}(\mathcal A)
\]
for the ordinary category of such models and enriched natural transformations.

\begin{proposition}[Models are subterminal objects]
\label{prop:models-sierpinski-sketch-subterminal}
For every categorical fuzzy frame \(\mathcal A\), evaluation at \(U\) induces an
equivalence of ordinary categories
\[
    \operatorname{Mod}_{\Omega}(\mathcal A)
    \simeq
    \operatorname{Sub}_{\mathcal A}(1_{\mathcal A}).
\]
\end{proposition}

\begin{proof}
Let
\[
    F:\mathbb T_{\Omega}\to\mathcal A
\]
be a model. Since the sketch declares \(U\) subterminal, \(F(U)\) is subterminal in
\(\mathcal A\). The formal map \(U\to 1\) is sent to a map
\[
    F(U)\to F(1)\cong 1_{\mathcal A}.
\]
Thus evaluation at \(U\) gives a functor
\[
    \operatorname{Mod}_{\Omega}(\mathcal A)
    \longrightarrow
    \operatorname{Sub}_{\mathcal A}(1_{\mathcal A}).
\]

Conversely, let
\[
    S\hookrightarrow 1_{\mathcal A}
\]
be a subterminal object of \(\mathcal A\). By the initiality of the free enriched
finite-limit-and-scalar sketch generated by one subterminal object, giving a model
\[
    F:\mathbb T_{\Omega}\to\mathcal A
\]
with \(F(U)=S\) is equivalent to interpreting the generator \(U\) as the chosen
subterminal object \(S\). The values on finite-limit formulae, scalar formulae, and
the generating maps \(\kappa_{V,A}\) are then forced up to the unique coherent
isomorphisms determined by the relevant universal properties.

Explicitly, the corresponding model \(F_S\) is defined recursively by
\[
    F_S(U)=S,
    \qquad
    F_S(1)=1_{\mathcal A},
\]
by interpreting the specified finite-limit formulae as the corresponding finite
limits in \(\mathcal A\), and by interpreting scalar formulae as tensors:
\[
    F_S(V\diamond A)
    =
    V\tensor_{\mathcal A}F_S(A).
\]
The generating morphism
\[
    \kappa_{V,A}:
    V\to \mathbb T_{\Omega}(A,V\diamond A)
\]
is interpreted as the tensor coevaluation
\[
    V\to
    \mathcal A(F_S(A),V\tensor_{\mathcal A}F_S(A)).
\]

Similarly, a morphism of models
\[
    F_S\to F_T
\]
is determined by its component at the generator \(U\). The components at
finite-limit formulae are forced by the universal property of the corresponding
limits, and the components at scalar formulae are forced by tensoring with the
scalar object. The component at \(U\) is exactly a morphism
\[
    S\to T
\]
over \(1_{\mathcal A}\). Therefore evaluation at \(U\) is fully faithful and
essentially surjective, hence an equivalence:
\[
    \operatorname{Mod}_{\Omega}(\mathcal A)
    \simeq
    \operatorname{Sub}_{\mathcal A}(1_{\mathcal A}).
\]
\end{proof}

\subsection{The Sierpiński localization}
\label{subsec:sierpinski-localization}

Form the fuzzy presheaf category
\[
    P_{\mathbb V}(\mathbb T_{\Omega})
    =
    [\mathbb T_{\Omega}^{\op},\underline{\mathbb V}].
\]
Let
\[
    y:\mathbb T_{\Omega}\to P_{\mathbb V}(\mathbb T_{\Omega})
\]
be the Yoneda embedding.

We use the standard enriched finite-limit sketch localization theorem in the
following form.

\begin{proposition}[Sierpiński sketch localization]
\label{prop:sierpinski-sketch-localization}
There is a left-exact reflective localization
\[
    a_{\Omega}\dashv i_{\Omega}:
    P_{\mathbb V}(\mathbb T_{\Omega})
    \rightleftarrows
    P_{\mathbb V}(\mathbb T_{\Omega})_{J_{\Omega}}
\]
with the following universal property. For every categorical fuzzy frame
\(\mathcal A\), restriction along the composite
\[
    \mathbb T_{\Omega}
    \xrightarrow{y}
    P_{\mathbb V}(\mathbb T_{\Omega})
    \xrightarrow{a_{\Omega}}
    P_{\mathbb V}(\mathbb T_{\Omega})_{J_{\Omega}}
\]
induces a natural equivalence of ordinary categories
\[
    \mathbf{FFrm}_{\mathbb V}
    \bigl(P_{\mathbb V}(\mathbb T_{\Omega})_{J_{\Omega}},\mathcal A\bigr)
    \simeq
    \operatorname{Mod}_{\Omega}(\mathcal A).
\]
\end{proposition}

\begin{proof}
This is the enriched finite-limit sketch localization theorem applied to the small
sketch \(\mathbb T_{\Omega}\). The localization is generated by the sketch data: the
specified finite-limit cones, the subterminality equation for \(U\), and the scalar
comparison maps
\[
    \alpha_{V,A}:
    V\tensor yA
    \longrightarrow
    y(V\diamond A).
\]
The universal property says exactly that small-colimit-preserving left-exact
\(\mathbb V\)-functors out of the localized presheaf category are the same as
models of the sketch in the target categorical fuzzy frame.
\end{proof}

The localization is generated by the enriched sketch data in the following sense:
after reflection, the designated finite-limit cones become limiting cones, the
subterminality equation for \(U\) holds, and the scalar comparison maps
\[
    \alpha_{V,A}:
    V\tensor yA
    \longrightarrow
    y(V\diamond A)
\]
become isomorphisms in the correct classifying variance. Thus the localization is
characterized by its universal property for small-colimit-preserving left-exact
functors out of the localized presheaf category.

It is important not to describe the \(J_{\Omega}\)-local presheaves themselves as
presheaves satisfying the tensor equations. Orthogonality of a presheaf \(F\) to
\[
    V\tensor yA\to y(V\diamond A)
\]
would instead identify
\[
    F(V\diamond A)
    \cong
    [V,F(A)].
\]
The tensor equations are imposed in the correct variance by the universal property
for cocontinuous left-exact functors out of the localization.

\begin{definition}[Fuzzy Sierpiński frame]
\label{def:fuzzy-sierpinski-frame}
The \textbf{fuzzy Sierpiński frame} over the fuzzy cosmos \(\mathbb V\) is
\[
    \Omega_{\mathbb V}
    \eqdef
    P_{\mathbb V}(\mathbb T_{\Omega})_{J_{\Omega}}.
\]
Its \textbf{generic fuzzy open} is
\[
    \sigma_{\mathbb V}
    \eqdef
    a_{\Omega}(yU).
\]
\end{definition}

By construction,
\[
    \Omega_{\mathbb V}
\]
is a left-exact reflective localization of a fuzzy presheaf category. Hence it is a
categorical fuzzy frame in the sense of this chapter.

\begin{proposition}[The generic fuzzy open is subterminal]
\label{prop:generic-fuzzy-open-subterminal}
The object
\[
    \sigma_{\mathbb V}\in\Omega_{\mathbb V}
\]
is subterminal.
\end{proposition}

\begin{proof}
The localization \(a_{\Omega}\dashv i_{\Omega}\) is generated so that the image of
the designated subterminality data for \(U\) becomes valid in
\(\Omega_{\mathbb V}\). Since the localization also forces the specified product cone
for \(U\times U\) to become a product cone, the object
\[
    a_{\Omega}y(U\times U)
\]
is canonically isomorphic to
\[
    \sigma_{\mathbb V}\times\sigma_{\mathbb V}.
\]
Moreover, the image of the diagonal comparison
\[
    U\to U\times U
\]
becomes an isomorphism after applying \(a_{\Omega}y\). Hence
\[
    \sigma_{\mathbb V}
    \cong
    \sigma_{\mathbb V}\times\sigma_{\mathbb V}.
\]
Since \(\Omega_{\mathbb V}\) has a terminal object, this says precisely that
\(\sigma_{\mathbb V}\) is subterminal.
\end{proof}

\subsection{The classifying theorem}
\label{subsec:sierpinski-classifying-theorem}

\begin{theorem}[Fuzzy Sierpiński classification, frame form]
\label{thm:fuzzy-sierpinski-classification-frame}
For every categorical fuzzy frame \(\mathcal A\), evaluation at the generic fuzzy
open induces a natural equivalence of ordinary categories
\[
    \mathbf{FFrm}_{\mathbb V}(\Omega_{\mathbb V},\mathcal A)
    \simeq
    \operatorname{Sub}_{\mathcal A}(1_{\mathcal A}).
\]
Under this equivalence, a fuzzy-frame morphism
\[
    H:\Omega_{\mathbb V}\to\mathcal A
\]
corresponds to the subterminal object
\[
    H(\sigma_{\mathbb V})\hookrightarrow 1_{\mathcal A}.
\]
\end{theorem}

\begin{proof}
By Proposition~\ref{prop:sierpinski-sketch-localization}, restriction along
\[
    a_{\Omega}y:\mathbb T_{\Omega}\to\Omega_{\mathbb V}
\]
induces a natural equivalence
\[
    \mathbf{FFrm}_{\mathbb V}(\Omega_{\mathbb V},\mathcal A)
    \simeq
    \operatorname{Mod}_{\Omega}(\mathcal A).
\]
Under this equivalence, a morphism
\[
    H:\Omega_{\mathbb V}\to\mathcal A
\]
corresponds to the model
\[
    Ha_{\Omega}y:\mathbb T_{\Omega}\to\mathcal A.
\]
Its value at the generator \(U\) is
\[
    H(a_{\Omega}yU)
    =
    H(\sigma_{\mathbb V}).
\]

By Proposition~\ref{prop:models-sierpinski-sketch-subterminal}, evaluation at \(U\)
gives an equivalence
\[
    \operatorname{Mod}_{\Omega}(\mathcal A)
    \simeq
    \operatorname{Sub}_{\mathcal A}(1_{\mathcal A}).
\]
Combining these equivalences gives
\[
    \mathbf{FFrm}_{\mathbb V}(\Omega_{\mathbb V},\mathcal A)
    \simeq
    \operatorname{Sub}_{\mathcal A}(1_{\mathcal A}).
\]
Under the composite equivalence, \(H\) is sent to
\[
    H(\sigma_{\mathbb V})\hookrightarrow 1_{\mathcal A}.
\]
Since \(H\) preserves terminal objects and binary products,
\(H(\sigma_{\mathbb V})\) is subterminal.
\end{proof}

\begin{definition}[Characteristic morphism of a fuzzy open]
\label{def:characteristic-morphism-fuzzy-open}
Let \(\mathcal A\) be a categorical fuzzy frame, and let
\[
    S\hookrightarrow 1_{\mathcal A}
\]
be a fuzzy open. The \textbf{characteristic morphism} of \(S\) is the fuzzy-frame
morphism
\[
    \chi_S^\ast:
    \Omega_{\mathbb V}\to\mathcal A
\]
corresponding to \(S\) under the equivalence of
Theorem~\ref{thm:fuzzy-sierpinski-classification-frame}. Equivalently,
\[
    \chi_S^\ast(\sigma_{\mathbb V})
    \cong
    S
\]
as subobjects of \(1_{\mathcal A}\).
\end{definition}

\subsection{The fuzzy Sierpiński locale}
\label{subsec:fuzzy-sierpinski-locale}

Define
\[
    \mathbf{FLoc}_{\mathbb V}
    \eqdef
    \mathbf{FFrm}_{\mathbb V}^{\op}.
\]

\begin{definition}[Fuzzy Sierpiński locale]
\label{def:fuzzy-sierpinski-locale}
The \textbf{fuzzy Sierpiński locale} over \(\mathbb V\) is the fuzzy locale
\[
    \mathbb S_{\mathbb V}
\]
whose fuzzy frame of opens is
\[
    \Opens_{\mathbb V}(\mathbb S_{\mathbb V})
    =
    \Omega_{\mathbb V}.
\]
\end{definition}

\begin{theorem}[Fuzzy Sierpiński classification, localic form]
\label{thm:fuzzy-sierpinski-classification-localic}
For every categorical fuzzy locale \(\mathcal X\), there is a natural equivalence
of ordinary categories
\[
    \mathbf{FLoc}_{\mathbb V}(\mathcal X,\mathbb S_{\mathbb V})
    \simeq
    \operatorname{Sub}_{\Opens_{\mathbb V}(\mathcal X)}
    \bigl(1_{\Opens_{\mathbb V}(\mathcal X)}\bigr).
\]
Thus maps into the fuzzy Sierpiński locale classify cosmos-level fuzzy opens.
\end{theorem}

\begin{proof}
By definition of fuzzy locales,
\[
\begin{aligned}
    \mathbf{FLoc}_{\mathbb V}(\mathcal X,\mathbb S_{\mathbb V})
    &=
    \mathbf{FFrm}_{\mathbb V}
    \bigl(
        \Opens_{\mathbb V}(\mathbb S_{\mathbb V}),
        \Opens_{\mathbb V}(\mathcal X)
    \bigr)
    \\
    &=
    \mathbf{FFrm}_{\mathbb V}
    \bigl(
        \Omega_{\mathbb V},
        \Opens_{\mathbb V}(\mathcal X)
    \bigr).
\end{aligned}
\]
Applying Theorem~\ref{thm:fuzzy-sierpinski-classification-frame} to
\[
    \mathcal A=\Opens_{\mathbb V}(\mathcal X)
\]
gives
\[
    \mathbf{FFrm}_{\mathbb V}
    \bigl(
        \Omega_{\mathbb V},
        \Opens_{\mathbb V}(\mathcal X)
    \bigr)
    \simeq
    \operatorname{Sub}_{\Opens_{\mathbb V}(\mathcal X)}
    \bigl(1_{\Opens_{\mathbb V}(\mathcal X)}\bigr).
\]
Therefore
\[
    \mathbf{FLoc}_{\mathbb V}(\mathcal X,\mathbb S_{\mathbb V})
    \simeq
    \operatorname{Sub}_{\Opens_{\mathbb V}(\mathcal X)}
    \bigl(1_{\Opens_{\mathbb V}(\mathcal X)}\bigr).
\]
\end{proof}

The characteristic map of a fuzzy open
\[
    S\hookrightarrow 1_{\Opens_{\mathbb V}(\mathcal X)}
\]
is the fuzzy-locale map
\[
    \chi_S:\mathcal X\to\mathbb S_{\mathbb V}
\]
whose inverse-image morphism
\[
    \chi_S^\ast:
    \Omega_{\mathbb V}\to\Opens_{\mathbb V}(\mathcal X)
\]
satisfies
\[
    \chi_S^\ast(\sigma_{\mathbb V})\cong S.
\]

\subsection{Relation with the raw two-point presheaf}
\label{subsec:relation-raw-two-point-presheaf}

The construction above explains why the fuzzy Sierpiński frame should not be
identified, in general, with the raw presheaf frame on the two-point subterminal
shape.

Let
\[
    C=(U\leq 1)
\]
be the ordinary two-element preorder, regarded as a \(\Lambda\)-preorder by taking
\[
    C(U,1)=\top_L,
    \qquad
    C(1,U)=\bot_L,
\]
and identity homs equal to \(\top_L\). In the crisp case, the presheaf frame
\[
    P_2(C)
\]
is the ordinary Sierpiński frame. For a general fuzzy truth-value base
\[
    \Lambda=(L,\leq,\meet,\join,\tensor,\multimap,\bot_L,\top_L),
\]
however, the raw presheaf frame
\[
    P_\Lambda(C)
\]
does not have the required free finite-meet behavior among scalar-generated fuzzy
open formulae.

In the posetal reflection, an element of
\[
    P_\Lambda(C)
\]
is a pair
\[
    (u,t)\in L\times L
\]
with
\[
    t\leq u,
\]
where \(u\) is the value at \(U\) and \(t\) is the value at \(1\). The representable
corresponding to \(U\) is
\[
    \sigma=(\top_L,\bot_L),
\]
and the top predicate is
\[
    \top=(\top_L,\top_L).
\]
Every element decomposes as
\[
    (u,t)
    =
    (u\tensor\sigma)\join(t\tensor\top).
\]
Indeed,
\[
    u\tensor\sigma=(u,\bot_L),
    \qquad
    t\tensor\top=(t,t),
\]
and \(t\leq u\).

Therefore any scalar-compatible frame homomorphism
\[
    H:P_\Lambda(C)\to A
\]
into a \(\Lambda\)-scalar frame \(A\), with
\[
    H(\sigma)=a,
\]
is forced to satisfy
\[
    H(u,t)
    =
    (u\tensor_A a)\join (t\tensor_A\top_A).
\]
In the special case \(A=\Lambda\), this becomes
\[
    H_a(u,t)=(u\tensor a)\join t.
\]
For non-idempotent tensors, this forced map need not preserve binary meets.

For example, let
\[
    \Lambda=[0,1]
\]
with the {\L}ukasiewicz tensor
\[
    r\tensor s=\max(0,r+s-1).
\]
Take
\[
    a=\frac12,
    \qquad
    p=(0.8,0.1),
    \qquad
    q=(0.4,0.3).
\]
Then
\[
    p\meet q=(0.4,0.1).
\]
The forced map
\[
    H_a(u,t)=(u\tensor a)\join t
\]
satisfies
\[
    H_a(p)=0.3,
    \qquad
    H_a(q)=0.3.
\]
Thus
\[
    H_a(p)\meet H_a(q)=0.3.
\]
But
\[
    H_a(p\meet q)
    =
    H_a(0.4,0.1)
    =
    0.1.
\]
Hence
\[
    H_a(p\meet q)\neq H_a(p)\meet H_a(q).
\]
So \(H_a\) is not a frame homomorphism.

Similarly, for the product tensor
\[
    r\tensor s=rs,
\]
take the same values
\[
    a=\frac12,
    \qquad
    p=(0.8,0.1),
    \qquad
    q=(0.4,0.3).
\]
Then
\[
    H_a(u,t)=\max(au,t).
\]
Thus
\[
    H_a(p)=\max(0.4,0.1)=0.4,
    \qquad
    H_a(q)=\max(0.2,0.3)=0.3,
\]
so
\[
    H_a(p)\meet H_a(q)=0.3.
\]
But
\[
    H_a(p\meet q)
    =
    H_a(0.4,0.1)
    =
    \max(0.2,0.1)
    =
    0.2.
\]
Hence
\[
    H_a(p\meet q)\neq H_a(p)\meet H_a(q).
\]
Therefore the raw two-point presheaf frame does not classify arbitrary fuzzy opens
for these non-idempotent bases.

The enriched syntactic site \(\mathbb T_{\Omega}\) corrects this defect by adding
formal scalar formulae and their finite-limit combinations before freely adjoining
presheaves and then imposing the sketch localization. Thus the fuzzy Sierpiński
frame is
\[
    \Omega_{\mathbb V}
    =
    P_{\mathbb V}(\mathbb T_{\Omega})_{J_{\Omega}},
\]
not the raw presheaf object on the two-point shape.

\begin{corollary}[Crisp specialization]
\label{cor:crisp-sierpinski-specialization}
When
\[
    \Lambda=2,
\]
the fuzzy Sierpiński frame \(\Omega_2\) is isomorphic to the ordinary three-element
Sierpiński frame. Equivalently, it is isomorphic to the downset frame of the
two-element preorder
\[
    U\leq 1.
\]
\end{corollary}

\begin{proof}
For \(\Lambda=2\), the scalar action is forced:
\[
    1\tensor a=a,
    \qquad
    0\tensor a=0.
\]
Thus scalar compatibility adds no structure beyond the ordinary finite-meet and
join structure. The free fuzzy frame on one fuzzy open is therefore the ordinary
free frame on one open, namely the Sierpiński frame. This is the downset frame of
the preorder \(U\leq 1\), whose generator is the open \(\{U\}\), and which has three
elements:
\[
    \varnothing,
    \qquad
    \{U\},
    \qquad
    \{U,1\}.
\]
\end{proof}

\subsection{Posetal reflection}
\label{subsec:sierpinski-posetal-reflection}

When the fuzzy cosmos is the thin cosmos associated to
\[
    \Lambda=(L,\leq,\meet,\join,\tensor,\multimap,\bot_L,\top_L),
\]
we write
\[
    \Omega_\Lambda,
    \qquad
    \sigma_\Lambda,
    \qquad
    \mathbb S_\Lambda
\]
for the corresponding posetal reflections of
\[
    \Omega_{\mathbb V},
    \qquad
    \sigma_{\mathbb V},
    \qquad
    \mathbb S_{\mathbb V}.
\]

The enriched syntactic sketch has a transparent posetal shadow. In the posetal
comparison where the only finite limits used explicitly are terminal objects and
binary meets, its formal objects are finite fuzzy-open formulae generated by
\[
    \top,
    \qquad
    \sigma,
    \qquad
    r\tensor\varphi \quad (r\in L),
    \qquad
    \varphi\meet\psi.
\]
Here \(\sigma\) denotes the formal generator corresponding to \(U\). Arbitrary joins
are not part of the finite syntactic sketch; they are freely adjoined by passing to
the presheaf frame and then imposing the sketch localization.

The fuzzy hom-value
\[
    \mathbb T_{\Omega}(\varphi,\psi)
\]
is the degree to which the sequent
\[
    \varphi\Rightarrow\psi
\]
is derivable from the finite-limit, subterminality, and scalar rules of the sketch.

The posetal fuzzy Sierpiński frame is
\[
    \Omega_\Lambda
    =
    P_\Lambda(\mathbb T_{\Omega})_{J_{\Omega}},
\]
with generic fuzzy open
\[
    \sigma_\Lambda
    =
    a_{\Omega}(y\sigma).
\]
For every fuzzy frame \(A\) over \(\Lambda\), the classification theorem specializes
to the set-level universal property
\[
    \FFrm_\Lambda(\Omega_\Lambda,A)
    \cong
    |A|.
\]
Here the right-hand side is the underlying set of the frame \(A\), because in the
posetal reflection every element \(a\in A\) determines a unique subobject
\[
    a\hookrightarrow \top_A,
\]
and every subobject of \(\top_A\) arises this way. The element corresponding to a
fuzzy-frame homomorphism
\[
    H:\Omega_\Lambda\to A
\]
is
\[
    H(\sigma_\Lambda)\in |A|.
\]

Dualizing, for every fuzzy locale \(\mathcal X\),
\[
    \FLoc_\Lambda(\mathcal X,\mathbb S_\Lambda)
    \cong
    |\Opens_\Lambda(\mathcal X)|.
\]
Thus maps into the fuzzy Sierpiński locale classify fuzzy opens exactly as in the
ordinary localic case, but the construction of the classifier uses the enriched
syntactic site rather than the naive two-point presheaf object.
